% concatenated from manuscript2.tex
\documentclass[reqno]{amsart}
%\usepackage{amsaddr}
% Elsevier layout
\usepackage[left=2.6cm,right=2.6cm,top=3cm,bottom=3cm]{geometry}
\usepackage[backref=page]{hyperref}
\usepackage[nameinlink,capitalize]{cleveref}
\usepackage{cite}
\usepackage{enumitem}
\usepackage{graphicx}
\usepackage{caption}
\usepackage{subcaption}
%\usepackage{refcheck} % for editing only
%\usepackage[pagewise]{lineno}\linenumbers % for review only
%\usepackage{color,soul} % highlighted for revision only

\theoremstyle{plain}
\newtheorem{theorem}{Theorem}
\newtheorem{proposition}[theorem]{Proposition}
\newtheorem{corollary}[theorem]{Corollary}
\newtheorem{lemma}[theorem]{Lemma}

\theoremstyle{definition}
\newtheorem{definition}{Definition}
\newtheorem{example}{Example}

\theoremstyle{remark}
\newtheorem{remark}{Remark}

%\numberwithin{equation}{section}
\allowdisplaybreaks


\begin{document}

\author{Phuoc Vinh Dinh}
\address{Phuoc Vinh Dinh$^{1,2,3}$ (ORCID: 0009-0000-6297-7798)\newline
	$^1$Faculty of Applied Sciences, Ho Chi Minh City University of Technology (HCMUT), 268 Ly Thuong Kiet Street, Dien Hong Ward, Ho Chi Minh City, Vietnam;\newline
	$^2$Vietnam National University, Ho Chi Minh City, Vietnam\newline
	$^3$Department of Mathematics, FPT University, Ho Chi Minh City, Vietnam}
\email{dpvinh.sdh251@hcmut.edu.vn, vinhdp2@fe.edu.vn}

\author{Phuong Le}
\address{Phuong Le$^{4,2}$ (ORCID: 0000-0003-4724-7118)\newline
	$^4$Faculty of Economic Mathematics, University of Economics and Law, Ho Chi Minh City, Vietnam; \newline
	$^2$Vietnam National University, Ho Chi Minh City, Vietnam}
\email{phuongl@uel.edu.vn}

\author{Tien Dung Nguyen}
\address{Tien Dung Nguyen$^{1,2}$\newline
	$^1$Faculty of Applied Sciences, Ho Chi Minh City University of Technology (HCMUT), 268 Ly Thuong Kiet Street, Dien Hong Ward, Ho Chi Minh City, Vietnam;\newline
	$^2$Vietnam National University, Ho Chi Minh City, Vietnam}
\email{dungnt@hcmut.edu.vn}

\title[A free boundary model for invasive and native species]{A free boundary model for invasive and native species under shifting climate in the weak competition case}

\date{\today}

\begin{abstract}
	We study a free boundary problem for a diffusive Lotka--Volterra competition system describing the invasion of a new species into the habitat of a native competitor, in a habitat that is shifted from unfavourable to favourable at a constant speed $c>0$ by climate change. Only the invader feels the shifting environment and only its range is governed by a Stefan-type free boundary, while the native species occupies the whole half line. We work throughout in the weak competition regime, in which the two species may coexist. We prove a spreading--vanishing dichotomy: either the invader spreads and the pair converges to the coexistence steady state $(u^*,v^*)$, or the invader vanishes and the native species recovers its carrying capacity. In the vanishing case we obtain the explicit bound $\lim_{t\to\infty}h(t)\le\frac{\pi}{2}\sqrt{d_1c_2/(a_1c_2-a_2c_1)}$, and we give criteria guaranteeing each alternative. When spreading occurs, we determine the exact asymptotic spreading speed: $\lim_{t\to\infty}h(t)/t=\min\{c,c_0\}$, where $c_0$ is the spreading speed of the corresponding homogeneous weak competition system. In particular the invasion is slowed down both by the competitor and by the climate shift, and the slower of the two mechanisms is the one that determines the speed. Numerical simulations illustrate the results.
\end{abstract}

\subjclass[2020]{35K51, 35R35, 35B40, 92D25, 35Q92}

\dedicatory{}

\keywords{free boundary problem, shifting environment, climate warming, weak competition, semi-wave, spreading--vanishing dichotomy, asymptotic spreading speed}

\maketitle
%\tableofcontents

\section{Introduction and statement of the results}\label{sec_1}

\subsection{Background}

Understanding how species respond to a changing climate is a central question in
spatial ecology, and reaction--diffusion equations have long been used to model
the shift of habitats and the propagation of populations; see
\cite{MR2471053,MR3691993,MR3377515} and the references therein. When a species
expands into new territory, the boundary of its range is not known in advance,
which naturally leads to free boundary problems. The starting point of this line
of research is the work of Du and Lin \cite{MR2607347}, where a diffusive
logistic model with a Stefan-type free boundary was introduced to describe the
spreading of a single species. They established a sharp spreading--vanishing
dichotomy: either the species spreads over the whole space and stabilises at its
carrying capacity, or it vanishes; which alternative occurs is governed by the
initial habitat size and by the expansion capacity of the species. We refer to
the survey \cite{MR4404216} for an account of the many developments that
followed.

To take the non-stationarity of the environment into account, several authors
have coupled the free boundary framework with a shifting habitat
\cite{MR3764580,MR3871607,MR3639148,MR4766631}. Of particular relevance to the
present paper is the model of Hu, Hao, Song and Du \cite{HU20205931},
\begin{equation}\label{single}
	\begin{cases}
		u_t = d_1 u_{xx} + (A(x - ct) - b_1 u) u, & t > 0,~ 0 < x < h(t), \\
		u_x(t, 0) = 0, ~ u(t, x) = 0, & t > 0, ~ h(t) \leq x < +\infty, \\
		h'(t) = -\mu (A(h(t) - ct))u_x(t, h(t)), & t > 0, \\
		h(0)=h_0,~ u(0, x) = u_0(x), & 0 \leq x \leq h_0,
	\end{cases}
\end{equation}
in which the environment turns from \emph{unfavourable to favourable} at speed
$c>0$, so that the species benefits from the climate shift, a situation which
is particularly relevant in invasion ecology, since many invasive species take
advantage of climate warming. For \eqref{single} a spreading--vanishing dichotomy
holds and, when spreading occurs, there is a critical speed
$\tilde c_0\in(0,2\sqrt{a_1d_1})$ such that the spreading profile is a semi-wave
with forced speed $c$ when $c<\tilde c_0$, and the usual semi-wave with speed
$\tilde c_0$ when $c\ge \tilde c_0$; in both cases the range expands at the
asymptotic speed $\min\{c,\tilde c_0\}$. (The notation $\tilde c_0$ is used here
only to avoid a clash with the speed $c_0$ of the competition system, introduced
in \eqref{c0} below.)

A second, independent extension of \cite{MR2607347} concerns the interaction
between an invader and a native competitor. In a homogeneous environment, Du and
Lin \cite{MR3327894} studied the two-species free boundary problem in which the
invader $u$ occupies the varying interval $[0,h(t)]$ while the resident $v$ lives
on the whole half line $[0,\infty)$. They showed that an inferior invader always
dies out, whereas a superior invader obeys a spreading--vanishing dichotomy; the
exact spreading speed of a superior invader was subsequently determined by Du,
Wang and Zhou \cite{MR3609207} through a semi-wave analysis, which appears to be
the first result of this kind for a system with a free boundary. The remaining
regime, \emph{weak competition}, in which the two species can coexist, was
treated by Wang, Nie and Du \cite{MR3986328}: a spreading--vanishing dichotomy
again holds, when the invasion succeeds both populations converge to the
coexistence state $(u^*,v^*)$, and the exact spreading speed exists and is
strictly smaller than the spreading speed of $u$ in the absence of $v$. Related
free boundary competition systems, including nonlocal ones and systems in which
both species have a free boundary, are studied in
\cite{Lei2018,MR4873231}.

The simultaneous effect of a shifting climate and of interspecific competition in
a free boundary setting was first addressed by Lei, Nie, Dong and Du
\cite{Lei2018}, where both competitors expand through free boundaries in a common
shifting habitat. In many real invasions, however, the native species is already
established across the whole habitat and does not experience a moving range
boundary; only the invader faces a climate-driven frontier. Motivated by this
observation, in paper \cite{DLN2026a} we studied
\begin{equation}\label{main}
	\begin{cases}
		u_t = d_1u_{xx} + (A(x - ct) - b_1u - c_1 v) u, & t > 0,~ 0 < x < h(t), \\
		v_t = d_2v_{xx} + (a_2 - b_2u - c_2v)v, & t > 0,~ 0 < x < +\infty, \\
		u_x(t, 0) = v_x(t, 0) = 0, ~ u(t, x) = 0, & t > 0, ~ h(t) \leq x < +\infty, \\
		h'(t) = -\mu (A(h(t) - ct))u_x(t, h(t)), & t > 0, \\
		h(0)=h_0,~ u(0, x) = u_0(x), & 0 \leq x \leq h_0, \\
		v(0, x) = v_0(x), & 0 \le x < +\infty,
	\end{cases}
\end{equation}
in the two \emph{weak--strong} competition regimes, that is, when the invader is
either an inferior or a superior competitor. The present paper is devoted to the
remaining case, and from the ecological point of view the most interesting,
case of \emph{weak competition}, in which neither species excludes the other and
coexistence is possible.

\subsection{The model and the standing assumptions}

In \eqref{main}, $u(t,x)$ and $v(t,x)$ denote the population densities of the
invasive and of the native species, and $x=h(t)$ is the free boundary, to be
determined together with $u$ and $v$. The invader exists initially on
$\{0\le x\le h_0\}$ and expands through the front $x=h(t)$, which obeys the
Stefan condition $h'(t)=-\mu(A(h(t)-ct))u_x(t,h(t))$; a derivation of this
condition from the consideration of population loss at the front may be found in
\cite{MR3764580}. The native species diffuses and grows on the whole habitat
$[0,\infty)$. Both species satisfy a no-flux condition at $x=0$. The constants
$d_1,d_2$ are the diffusion rates, $a_2$ is the intrinsic growth rate of $v$,
$b_1,c_2$ are the intra-specific and $c_1,b_2$ the inter-specific competition
rates, all assumed positive; $h_0>0$.

The intrinsic growth rate of the invader is $A(x-ct)$: the environment is shifted
at the constant speed $c>0$ in the direction of increasing $x$. Here $A$ should be
read as describing the \emph{climatic range limit of the invader}, which is what
moves, rather than climate change acting on the community as a whole: the native
competitor is already established throughout $[0,\infty)$ and, on the scales
relevant to the invasion, sits well inside its own climatic envelope, so that its
growth rate is effectively constant. See \cite[Remark 4]{DLN2026a} for a fuller
discussion, and \cite{Lei2018} for the complementary situation in which both
species are subject to a shifting habitat. Throughout the
paper we assume, as in \cite{HU20205931,DLN2026a}, that $A$ is Lipschitz
continuous on $\mathbb{R}$, strictly decreasing on $[0,l_0]$, and
\begin{equation}\label{A}
	A(\xi) =
	\begin{cases}
		a_1 & \text{ if } \xi \leq 0,\\
		a_0 & \text{ if } \xi \geq l_0,
	\end{cases}
\end{equation}
where $l_0 > 0$ and $a_0 < 0 < a_1$. Thus the region $\{x<ct\}$ is favourable for
the invader and the region $\{x>ct+l_0\}$ is unfavourable, and the favourable
region advances at speed $c$. The function $\mu$, which measures how easily the
range boundary expands, is Lipschitz continuous and increasing on $[a_0,a_1]$
with
\begin{equation}\label{mu}
	0 < \mu(a_0) \le \mu(a_1).
\end{equation}
It is convenient (and harmless) to extend $\mu$ to all of $\mathbb{R}$ by setting
$\mu(s):=\mu(a_0)$ for $s<a_0$ and $\mu(s):=\mu(a_1)$ for $s>a_1$; the extended
function is again Lipschitz and nondecreasing, and this extension is used only in
the construction of the perturbed semi-waves in \cref{sec_5}.

The initial functions satisfy
\begin{equation}\label{initial}
	\begin{cases}
		u_0 \in C^2([0,h_0]),& u_0 > 0 \text{ in } [0,h_0), ~ u_0'(0) = u_0(h_0) = 0,\\
		v_0 \in C^2([0,+\infty)) \cap L^\infty([0,+\infty)),& v_0 > 0 \text{ in } [0,+\infty), ~ v_0'(0) = 0,\\
		&\displaystyle\liminf_{x\to+\infty} v_0(x) > 0.
	\end{cases}
\end{equation}
The last requirement in \eqref{initial} says that the native species is already
established over the whole habitat; together with the continuity and the
positivity of $v_0$ it yields
\begin{equation}\label{v0inf}
	\sigma_0^{\max}:=\inf_{x\ge0} v_0(x) > 0 .
\end{equation}
This hypothesis is the exact analogue of \cite[(1.4)]{MR3986328} and of
\cite[(1.6)]{MR3609207}; it is used only in \cref{sec_6}, where the spreading
speed is determined, and it cannot be dispensed with there (see
\cite[Remark 3.5]{MR3609207}).

When the native competitor is absent ($v\equiv0$), \eqref{main} reduces to
\eqref{single}, studied in \cite{HU20205931}. When the shifting term is replaced
by the constant $a_1$, \eqref{main} reduces to the homogeneous competition system
of Du and Lin \cite{MR3327894},
\begin{equation}\label{main2}
	\begin{cases}
		u_t = d_1u_{xx} + (a_1 - b_1u - c_1 v) u, & t > 0,~ 0 < x < h(t), \\
		v_t = d_2v_{xx} + (a_2 - b_2u - c_2v)v, & t > 0,~ 0 < x < +\infty, \\
		u_x(t, 0) = v_x(t, 0) = 0, ~ u(t, x) = 0, & t > 0, ~ h(t) \leq x < +\infty, \\
		h'(t) = -\mu (a_1)u_x(t, h(t)), & t > 0, \\
		h(0)=h_0,~ u(0, x) = u_0(x), & 0 \leq x \leq h_0, \\
		v(0, x) = v_0(x), & 0 \le x < +\infty.
	\end{cases}
\end{equation}
Problem \eqref{main2} was analysed in \cite{MR3327894,MR3609207} under the
weak--strong competition conditions
$\frac{a_1}{a_2} < \min\{\frac{b_1}{b_2},\frac{c_1}{c_2}\}$ (inferior invader) or
$\frac{a_1}{a_2} > \max\{\frac{b_1}{b_2},\frac{c_1}{c_2}\}$ (superior invader),
and in \cite{MR3986328} under the weak competition condition
\begin{equation}\label{weak_competition}
	\frac{c_1}{c_2} < \frac{a_1}{a_2} < \frac{b_1}{b_2},
\end{equation}
which we assume from now on. Under \eqref{weak_competition} the kinetic system
associated with \eqref{main2} has the globally attracting coexistence state
\begin{equation}\label{coexist}
	u^* := \frac{a_1c_2-a_2c_1}{b_1c_2-b_2c_1},\qquad
	v^* := \frac{a_2b_1-a_1b_2}{b_1c_2-b_2c_1},
\end{equation}
and \eqref{weak_competition} guarantees $b_1c_2-b_2c_1>0$, $u^*>0$, $v^*>0$ and
\begin{equation}\label{ineq_star}
	0<u^*<\frac{a_1}{b_1},\qquad 0<v^*<\frac{a_2}{c_2},\qquad
	a_1-\frac{a_2c_1}{c_2}>0 .
\end{equation}
Throughout the paper we write
\begin{equation}\label{Rstar}
	R^*:=\frac{\pi}{2}\sqrt{\frac{d_1c_2}{a_1c_2-a_2c_1}},
	\qquad
	R_0:=\frac{\pi}{2}\sqrt{\frac{d_1}{a_1}},
\end{equation}
so that $0<R_0<R^*$ by \eqref{ineq_star}.

\subsection{Main results}

Our first result describes the long-time dynamics of \eqref{main}. Recall from
\cref{th6} below that $h$ is strictly increasing, so that
$h_\infty:=\lim_{t\to\infty}h(t)\in(h_0,+\infty]$ always exists.

\begin{theorem}\label{th:dynamics}
	Let $(u, v, h)$ be the unique solution of \eqref{main} with initial functions
	satisfying \eqref{initial}, and suppose that \eqref{weak_competition} holds.
	Then exactly one of the following alternatives occurs.
	\begin{enumerate}
		\item[(i)] \textit{Spreading of $u$}:
		\[
		h_\infty = +\infty \quad\text{ and }\quad \lim_{t\to\infty} (u(t,\cdot), v(t,\cdot)) = (u^*, v^*) \text{ in } C^2_{\rm loc}([0, \infty));
		\]
		\item[(ii)] \textit{Vanishing of $u$}:
		\[
		h_\infty \le R^* \quad\text{ and }\quad
		\lim_{t\to\infty} \|u(t,\cdot)\|_{C([0,h(t)])}=0,\quad
		\lim_{t\to\infty} v(t,\cdot) = \frac{a_2}{c_2} \text{ in } C^2_{\rm loc}([0, \infty)).
		\]
	\end{enumerate}
\end{theorem}

The next result gives criteria under which each of the two alternatives takes
place. It is the analogue, for the competition system \eqref{main}, of
\cite[Theorem 1.5]{HU20205931} and of \cite[Theorem 1.2]{MR3986328}.

\begin{theorem}\label{th:criteria}
	Assume \eqref{weak_competition} and let $(u,v,h)$ be the solution of
	\eqref{main} with initial functions satisfying \eqref{initial}.
	\begin{enumerate}
		\item[(i)] If $h_0\ge R^*$, then spreading always happens.
		\item[(ii)] Fix $h_0<R_0$ and $v_0$, write $u_0=\sigma\phi$ where $\phi$
		satisfies the conditions imposed on $u_0$ in \eqref{initial} and
		$\sigma>0$, and denote by $(u^\sigma,v^\sigma,h^\sigma)$ the corresponding
		solution. Then there exists
		$\sigma_0=\sigma_0(h_0,\phi,v_0,c)\in(0,+\infty]$ such that vanishing
		happens for $0<\sigma\le\sigma_0$ and spreading happens for
		$\sigma>\sigma_0$.
	\end{enumerate}
\end{theorem}

\begin{remark}\label{rm:criteria}
	As in \cite[Remark 1.6]{HU20205931} we do not know whether
	$\sigma_0=+\infty$ can actually occur; this remains an open problem already
	for the single species model \eqref{single}.
\end{remark}

To state our second main result we recall from \cite[Theorem 1.5]{MR3986328}
(see \cref{th:semi-wave0} below) that, under \eqref{weak_competition}, the
homogeneous problem \eqref{main2} possesses a well-defined asymptotic spreading
speed $c_0>0$, characterised by the semi-wave identity
$-\mu(a_1)\Psi_{c_0}'(0)=c_0$.

\begin{theorem}\label{th:spreading_speed}
	Assume \eqref{weak_competition} and \eqref{initial}. If alternative (i) of
	\cref{th:dynamics} occurs, then
	\[
	\lim_{t\to\infty} \frac{h(t)}{t} = \min\{c, c_0\}.
	\]
\end{theorem}

Let $c_0^{\rm sing}$ denote the asymptotic spreading speed of the invader in a
globally favourable environment \emph{in the absence of the native species},
that is, the speed at which the free boundary of
\begin{equation}\label{single_hom}
	\begin{cases}
		\hat u_t = d_1\hat u_{xx}+(a_1-b_1\hat u)\hat u, & t>0,\ 0<x<\hat h(t),\\
		\hat u_x(t,0)=\hat u(t,\hat h(t))=0,\quad \hat h'(t)=-\mu(a_1)\hat u_x(t,\hat h(t)), & t>0,\\
		\hat h(0)=h_0,\quad \hat u(0,\cdot)=u_0
	\end{cases}
\end{equation}
advances when spreading occurs; see \cite{MR2607347,MR4404216}.

\begin{lemma}\label{lm:c0sing}
	Assume \eqref{weak_competition}. Then $c_0\le c_0^{\rm sing}$.
\end{lemma}

\begin{proof}
	Let $(\tilde u,\tilde v,\tilde h)$ solve \eqref{main2} and let
	$(\hat u,\hat h)$ solve \eqref{single_hom} with the same $u_0$ and $h_0$.
	Since $\tilde v>0$ we have
	$\tilde u_t\le d_1\tilde u_{xx}+(a_1-b_1\tilde u)\tilde u$, while
	$\tilde u_x(t,0)=0$ and $\tilde h'(t)=-\mu(a_1)\tilde u_x(t,\tilde h(t))$;
	thus $(\tilde u,\tilde h)$ is a lower solution of \eqref{single_hom} and the
	scalar comparison principle gives $\tilde h\le\hat h$ for all $t>0$. If
	spreading occurs for \eqref{main2} then $\hat h_\infty=+\infty$ as well;
	dividing by $t$ and letting $t\to\infty$ gives $c_0\le c_0^{\rm sing}$.
\end{proof}

\begin{remark}\label{rm:speed}
	The formula $\min\{c,c_0\}$ has a transparent ecological meaning. The number
	$c_0$ is the speed at which the invader would advance in a globally
	favourable environment while competing with the native species, whereas $c$
	is the speed at which the favourable habitat itself moves. If $c_0\le c$ the
	invasion front travels at its own intrinsic speed and does not feel the
	climate constraint; if $c_0>c$ the front cannot outrun the moving edge of the
	favourable region and is slowed down to the speed $c$ of the climate shift.
	Note moreover that, by \cref{lm:c0sing}, $c_0\le c_0^{\rm sing}$, and that the
	inequality is in general strict: the native species depresses the invader
	precisely at its leading edge, where $v$ is close to its carrying capacity
	$a_2/c_2$. For the parameters used in \cref{sec_7} one finds
	$c_0\approx0.0128$ against $c_0^{\rm sing}\approx0.040$, a factor of more than
	three. Hence in the present model the invasion is retarded by two independent
	mechanisms, interspecific competition and the pace of climate change, and
	it is the more restrictive of the two that determines the observed speed.
\end{remark}

\begin{remark}\label{rm:compare}
	It is instructive to compare \cref{th:dynamics} with the corresponding result
	of \cite{DLN2026a} for a superior invader. There, when spreading occurs, the
	native species is driven to extinction and $(u,v)\to(a_1/b_1,0)$, whereas here
	the two species coexist and $(u,v)\to(u^*,v^*)$. This difference is not
	merely cosmetic: it changes the structure of the semi-waves that govern the
	front. In the superior case the $v$-component of the semi-wave vanishes at
	$-\infty$, and the construction of a lower solution requires a delicate
	modification of the profile far behind the front (see
	\cite[Lemma 3.3]{MR3609207} and the corresponding lemma in \cite{DLN2026a}). In
	the weak competition case the $v$-component tends to $v^*>0$ at $-\infty$,
	so no such modification is needed; on the other hand the state at $-\infty$ is
	now a genuine two-species coexistence equilibrium, and the asymptotic analysis
	of the semi-wave near $-\infty$ involves the full four-dimensional linearised
	system rather than a scalar one (see \cref{lm:asym}). A further difference
	appears in the vanishing case: since $u$ no longer eliminates $v$, the limit
	of $v$ is its own carrying capacity $a_2/c_2$ and not $0$.
\end{remark}

\begin{remark}\label{rm:hinfty}
	The bound $h_\infty\le R^*$ in \cref{th:dynamics}(ii) is sharper than the mere
	finiteness of $h_\infty$ and is what makes \cref{th:criteria} and the
	borderline case $c=c_0$ of \cref{th:spreading_speed} accessible; see
	\cref{lm:eq}. It is obtained from the following simple but useful observation
	(\cref{lem:reduction}): if $h_\infty<+\infty$, then $h(t)<ct$ for all large
	$t$, so that after a finite time the shifting environment is no longer felt
	and \eqref{main} reduces \emph{exactly} to the homogeneous system
	\eqref{main2}. The same observation is the key to the case $c>c_0$.
\end{remark}

\textbf{Organisation of the paper.} \Cref{sec_2} collects the basic results on
\eqref{main} (existence, uniqueness, global bounds) and proves, in full, the
comparison principle in the form used later; the details of the local existence
theory, which follow \cite[Theorem 2.1]{MR3327894}, are deferred to
\cref{sec_app}. \Cref{sec_3} contains the proof
of \cref{th:dynamics}, and \cref{sec_4} that of \cref{th:criteria}. \Cref{sec_5}
is devoted to the semi-waves associated with the shifting environment; this is
the technical core of the paper. \Cref{sec_6} proves \cref{th:spreading_speed}.
\Cref{sec_7} presents numerical simulations, and \cref{sec_8} concludes with some
comments and open questions.

\section{Preliminary results}\label{sec_2}

\subsection{Existence, uniqueness and global bounds}

\begin{theorem}\label{th4}
	For any given $\alpha \in (0, 1)$ and any $(u_0, v_0)$ satisfying
	\eqref{initial}, there exists a $T > 0$ such that problem \eqref{main} admits
	a unique bounded solution
	$$
	(u, v, h) \in C^{(1+\alpha)/2, 1 + \alpha}(D_T) \times C^{(1+\alpha)/2, 1 + \alpha}(D_T^{\infty}) \times C^{1 + \alpha/2}([0,T]).
	$$
	Moreover,
	$$
	\|u\|_{C^{(1+\alpha)/2, 1 + \alpha}(D_T)} + \|v\|_{C^{(1+\alpha)/2, 1 + \alpha}(D_T^{\infty})} + \|h\|_{C^{1 + \alpha/2}([0,T])} \leq C,
	$$
	where $D_T = \{(t,x) \in \mathbb{R}^2: t \in [0,T], x \in [0, h(t)]\}$,
	$D_T^{\infty} = \{(t,x)\in \mathbb{R}^2: t \in [0,T], x \in [0, +\infty)\}$, and $C$
	and $T$ depend only on $\alpha, h_0, c, \|u_0\|_{C^2([0,h_0])},
	\|v_0\|_{C^2([0,\infty))}$, on $a_0,a_1,a_2,b_1,b_2,c_1,c_2,d_1,d_2$ and on
	the Lipschitz constants of $A$ and $\mu$. Here $C^{(1+\alpha)/2,1+\alpha}$
	denotes the parabolic H\"older space of functions that are H\"older continuous
	with exponent $\frac{1+\alpha}{2}$ in $t$ and $1+\alpha$ in $x$, in the sense
	of \cite[Chapter IV]{MR241821}.
\end{theorem}

The proof is obtained by straightening the free boundary and applying the
contraction mapping theorem, exactly as in \cite[Theorem 2.1]{MR3327894}; the
only modification is the presence of the Lipschitz functions $A(\cdot-ct)$ and
$\mu(A(\cdot))$, which affects the constants but not the structure of the
argument. For the reader's convenience the details are given in \cref{sec_app}.

\begin{theorem}\label{th6}
	Problem \eqref{main} admits a unique solution $(u,v,h)$, defined for all
	$t>0$, and there exist positive constants $M_1,M_2,M_3$ such that
	\[
	0 < u(t,x) \leq M_1 \ \text{ for } (t, x) \in (0, +\infty) \times [0, h(t)),\qquad
	0 < v(t,x) \leq M_2 \ \text{ for } (t, x) \in (0, +\infty) \times [0, +\infty),
	\]
	\[
	0 < h'(t) \leq M_3 \ \text{ for } t \in (0,+\infty).
	\]
	Moreover $M_1=\max\{\|u_0\|_\infty,a_1/b_1\}$ and
	$M_2=\max\{\|v_0\|_\infty,a_2/c_2\}$ may be taken, and \eqref{main} has no
	unbounded solution.
\end{theorem}

\begin{proof}
	Let $[0,T_{\max})$ be the maximal interval of existence given by \cref{th4}.
	Since $A(\cdot)\le a_1$, the function $u$ satisfies
	$u_t-d_1u_{xx}\le u(a_1-b_1u)$ in $\{0<x<h(t)\}$ together with
	$u_x(t,0)=0$ and $u(t,h(t))=0$; comparing $u$ with the solution
	$\overline{u}(t)$ of the logistic equation
	\[
	\frac{{\rm d}\overline{u}}{{\rm d}t} = \overline{u}(a_1 - b_1\overline{u}),
	\qquad \overline{u}(0)=\|u_0\|_\infty,
	\qquad\text{i.e.}\qquad
	\overline{u}(t)=\frac{a_1e^{a_1t}}{b_1\left(e^{a_1t}-1+\frac{a_1}{b_1\|u_0\|_\infty}\right)},
	\]
	we obtain $u\le \sup_{t\ge0}\overline{u}(t)=M_1$. Similarly, $v$ satisfies
	$v_t-d_2v_{xx}\le v(a_2-c_2v)$ on $[0,\infty)$, whence $v\le M_2$. The strong
	maximum principle and the Hopf boundary lemma give $u>0$ in
	$\{0\le x<h(t)\}$, $u_x(t,h(t))<0$ and hence $h'(t)>0$, and also $v>0$.

	To bound $h'$ from above, fix $M\ge\max\{\sqrt{(a_1+c_1M_2)/(2d_1)},\,
	4\|u_0\|_{C^1([0,h_0])}/(3M_1)\}$ and set
	$\Omega:=\{(t,x):0<t<T_{\max},\,h(t)-M^{-1}<x<h(t)\}$ and
	$\overline{U}(t,x):=M_1[2M(h(t)-x)-M^2(h(t)-x)^2]$. A direct computation using
	$h'>0$ shows that $\overline{U}_t-d_1\overline{U}_{xx}\ge2d_1M_1M^2\ge
	(A(x-ct)-b_1u-c_1v)u$ in $\Omega$, while
	$\overline{U}(t,h(t)-M^{-1})=M_1\ge u$, $\overline{U}(t,h(t))=0=u(t,h(t))$ and
	$\overline{U}(0,\cdot)\ge u_0$ on $[h_0-M^{-1},h_0]$. The maximum principle
	yields $u\le\overline{U}$ in $\Omega$, hence
	$u_x(t,h(t))\ge\overline{U}_x(t,h(t))=-2MM_1$ and
	\[
	0<h'(t)=-\mu(A(h(t)-ct))u_x(t,h(t))\le 2MM_1\mu(a_1)=:M_3 .
	\]
	Since $M_1,M_2,M_3$ do not depend on $T_{\max}$, standard parabolic
	regularity and \cref{th4} allow the solution to be continued beyond
	$T_{\max}$ unless $T_{\max}=\infty$; therefore $T_{\max}=\infty$. Finally,
	comparing $u$ (extended by $0$ for $x>h(t)$) with $\|u_0\|_\infty+U(\cdot-x_0)$,
	where $U$ is the unique boundary blow-up solution of $-d_1U''=U(a_1-b_1U)$ in
	$(-1,1)$ with $U=+\infty$ on $\{\pm1\}$ (see \cite{DU_MA_2001}), and similarly
	for $v$, shows that every solution of \eqref{main} is automatically bounded.
\end{proof}

\subsection{Comparison principle}

The following comparison principle is the two-species analogue of
\cite[Lemma 2.6]{MR3327894}; note the \emph{interleaved} structure, typical of
competitive systems: an upper solution for $u$ is paired with a lower solution
for $v$, and vice versa. The two alternative conditions at $x=0$ are both needed:
the first is used for the upper solutions of \cref{sec_6}, the second for the
lower solutions, whose $u$-component is decreasing in $x$ and therefore does not
satisfy a one-sided Neumann condition at the origin.

\begin{theorem}\label{th7}
	Let $T \in (0,+\infty)$, and let
	\begin{align*}
		&\underline{h}, \overline{h} \in C^1([0,T]),\quad \underline h(0)>0,\\
		&\underline{u} \in C(\overline{U_T}) \cap C^{1,2}(U_T) \text{ with } U_T := \{ (t,x) \in \mathbb{R}^2 : t \in (0, T], x \in (0, \underline{h}(t))\},\\
		&\overline{u} \in C(\overline{V_T}) \cap C^{1,2}(V_T) \text{ with } V_T := \{ (t,x) \in \mathbb{R}^2 : t \in (0, T], x \in (0, \overline{h}(t))\},\\
		&\underline{v}, \overline{v} \in (L^{\infty} \cap C)([0,T] \times [0,\infty)) \cap C^{1,2}((0,T] \times [0,\infty)),\\
		&\underline{v}, \overline{v}, \underline{u}, \overline{u} \ge 0 .
	\end{align*}
	Let $(u,v,h)$ be the unique solution of \eqref{main}.

	\emph{(a)} Assume that
	\[
	\begin{cases}
		\overline{u}_t \ge d_1\overline{u}_{xx} + (A(x - ct) - b_1\overline{u} - c_1 \underline{v}) \overline{u}, & 0 < t \le T,~ 0 < x < \overline{h}(t), \\
		\underline{v}_t \le d_2\underline{v}_{xx} + (a_2 - b_2\overline{u} - c_2\underline{v})\underline{v}, & 0 < t \le T,~ 0 < x < +\infty, \\
		\overline{u}(t, x) = 0, & 0 < t \le T, ~ \overline{h}(t) \leq x < +\infty, \\
		\overline{h}'(t) \ge -\mu (A(\overline{h}(t) - ct))\overline{u}_x(t, \overline{h}(t)), & 0 < t \le T, \\
		\overline{h}(0) \ge h_0,~ \overline{u}(0, x) \ge u_0(x), & 0 \leq x \leq h_0, \\
		\underline{v}(0, x) \le v_0(x), & 0 \le x < +\infty,
	\end{cases}
	\]
	and that, in addition, either
	\[
	\overline{u}_x(t, 0) \le 0 \le \underline{v}_x(t, 0)\quad (0<t\le T)
	\qquad\text{or}\qquad
	\overline{u}(t,0)\ge u(t,0),\ \ \underline v(t,0)\le v(t,0)\quad (0<t\le T).
	\]
	Then
	\[
	h(t) \le \overline{h}(t) \text{ in } (0,T],\quad u(t,x) \le \overline{u} (t,x), ~ v(t,x) \ge \underline{v} (t,x) \text{ for } (t,x) \in (0,T] \times [0,+\infty).
	\]

	\emph{(b)} Assume that
	\[
	\begin{cases}
		\underline{u}_t \le d_1\underline{u}_{xx} + (A(x - ct) - b_1\underline{u} - c_1 \overline{v}) \underline{u}, & 0 < t \le T,~ 0 < x < \underline{h}(t), \\
		\overline{v}_t \ge d_2\overline{v}_{xx} + (a_2 - b_2\underline{u} - c_2\overline{v})\overline{v}, & 0 < t \le T,~ 0 < x < +\infty, \\
		\underline{u}(t, x) = 0, & 0 < t \le T, ~ \underline{h}(t) \leq x < +\infty, \\
		\underline{h}'(t) \le -\mu (A(\underline{h}(t) - ct))\underline{u}_x(t, \underline{h}(t)), & 0 < t \le T, \\
		\underline{h}(0) \le h_0,~ \underline{u}(0, x) \le u_0(x), & 0 \leq x \leq \underline h(0), \\
		\overline{v}(0, x) \ge v_0(x), & 0 \le x < +\infty,
	\end{cases}
	\]
	and that, in addition, either
	\[
	\underline{u}_x(t, 0) \ge 0 \ge \overline{v}_x(t, 0)\quad (0<t\le T)
	\qquad\text{or}\qquad
	\underline{u}(t,0)\le u(t,0),\ \ \overline v(t,0)\ge v(t,0)\quad (0<t\le T).
	\]
	Then
	\[
	h(t) \ge \underline{h}(t) \text{ in } (0,T],\quad u(t,x) \ge \underline{u} (t,x), ~ v(t,x) \le \overline{v} (t,x) \text{ for } (t,x) \in (0,T] \times [0,+\infty).
	\]
\end{theorem}

\begin{proof}
	We prove (a); part (b) is obtained by exchanging the roles of the two triples.
	Let $\tilde{M}$ be a common upper bound of $v$ and $\underline{v}$ on
	$[0,T]\times[0,+\infty)$ and set $w:=\tilde{M}-v$, $\overline{w} :=
	\tilde{M}-\underline{v}$. Since $s\mapsto \tilde M-s$ reverses the order and
	the resulting system for $(u,w)$ is cooperative, it suffices to prove that
	$u\le\overline u$ and $w\le\overline w$.

	\textit{Step 1: the case $h_0<\overline h(0)$.} We claim that
	$h(t)<\overline{h}(t)$ for all $t\in(0,T]$. If not, there is a first
	$t^*\in(0,T]$ with $h(t)<\overline{h}(t)$ on $(0,t^*)$ and
	$h(t^*)=\overline{h}(t^*)$; consequently
	\begin{equation}\label{eq:h}
		h'(t^*) \ge \overline{h}'(t^*).
	\end{equation}
	Set $U := (\overline{u} - u)e^{-Kt}$ and $W := (\overline{w} - w)e^{-Kt}$. By
	the mean value theorem,
	\begin{equation}\label{eq:sup2}
		\begin{cases}
			U_t - d_1 U_{xx} \ge -KU + c_{11}U + c_{12}W, & 0 < t \le t^*, ~ 0 \le x < h(t), \\
			W_t - d_2 W_{xx} \ge -KW + c_{21}U + c_{22}W, & 0 < t \le t^*, ~ 0 \le x < \infty, \\
			U(t,x) \ge 0, & 0 < t \le t^*, ~ h(t) \le x < \infty, \\
			U(0,x) \ge 0, ~ W(0,x) \ge 0, & 0 \le x < \infty,
		\end{cases}
	\end{equation}
	where
	\begin{align*}
		c_{11} &= A(x-ct) - b_1 (u+\overline u) - c_1 (\tilde{M} - \eta_1), & c_{12} &= c_1 \xi_1, \\
		c_{21} &= b_2 (\tilde{M} - \eta_2), & c_{22} &= a_2 - b_2 \xi_2 - c_2\big((\tilde M-\eta_2)+(\tilde M-\eta_2')\big),
	\end{align*}
	with $\xi_i$ between $u$ and $\overline{u}$ and $\tilde{M}-\eta_i$,
	$\tilde M-\eta_i'$ between $v$ and $\underline{v}$. All the $c_{ij}$ are
	bounded on $[0,t^*]\times[0,\infty)$ by \cref{th6} and the boundedness of the
	comparison functions, and, crucially,
	\begin{equation}\label{eq:coop}
		c_{12}\ge0,\qquad c_{21}\ge0 .
	\end{equation}
	We fix $K\ge 1+|c_{11}|+|c_{12}|+|c_{21}|+|c_{22}|$ on
	$[0,t^*]\times[0,\infty)$. Note that the third line of \eqref{eq:sup2} holds
	because $u\equiv0$ and $\overline u\ge0$ for $x\ge h(t)$; that $\overline u\ge0$
	follows from the maximum principle applied on
	$\{0\le x\le\overline h(t),\,0\le t\le T\}$.

	For $\ell>0$ define, with $d:=\max\{d_1,d_2\}$,
	\[
	\overline{U}(t,x) = U(t,x) + \frac{\tilde{M}(x^2 + 2dt)}{\ell^2},
	\qquad
	\overline{W}(t,x) = W(t,x) + \frac{\tilde{M}(x^2 + 2dt)}{\ell^2}.
	\]
	Using $\partial_t\frac{\tilde M(x^2+2dt)}{\ell^2}-d_i\partial_{xx}
	\frac{\tilde M(x^2+2dt)}{\ell^2}=\frac{2\tilde M(d-d_i)}{\ell^2}\ge0$,
	the pair $(\overline{U},\overline{W})$ satisfies on $[0,t^*]\times[0,\ell]$
	the same differential inequalities as $(U,W)$ up to a nonnegative right hand
	side, together with
	\[
	\overline{U}(t,x) \ge \tfrac{\tilde{M}(x^2 + 2dt)}{\ell^2} > 0 \ \ (h(t) \le x \le \ell),\qquad
	\overline{W}(t,\ell) \ge W(t,\ell) + \tilde{M} > 0,\qquad
	\overline{U}(0,\cdot),\overline{W}(0,\cdot)\ge0 ,
	\]
	the last inequality on $\{x=\ell\}$ because $|W|\le\tilde M$. Moreover, in the
	first alternative at $x=0$ we have
	$\overline U_x(t,0)\le0$ and $\overline W_x(t,0)\le0$, and in the second we
	have $\overline U(t,0)\ge0$ and $\overline W(t,0)\ge0$; in either case a
	negative minimum of $\overline U$ or $\overline W$ over
	$[0,t^*]\times[0,\ell]$ cannot be attained at $x=0$.

	Set
	$\kappa := \min\{\min_{[0,t^*]\times[0,\ell]}\overline{U},\,
	\min_{[0,t^*]\times[0,\ell]}\overline{W}\}$ and suppose $\kappa<0$. By the
	above, the minimum is attained either at some $(t_0,x_0)$ with
	$0<t_0\le t^*$, $0\le x_0<h(t_0)$ and $\overline U(t_0,x_0)=\kappa$, or at some
	$(t_1,x_1)$ with $0<t_1\le t^*$, $0\le x_1<\ell$ and
	$\overline W(t_1,x_1)=\kappa$. In the first case
	$\overline{U}_t(t_0,x_0)-d_1\overline{U}_{xx}(t_0,x_0)\le0$, whereas by
	\eqref{eq:coop} and the choice of $K$,
	\[
	-K\overline{U} + c_{11} \overline{U} + c_{12}\overline{W}
	\ \ge\ (-K + |c_{11}|)\kappa + c_{12}\kappa \ \ge\ -\kappa > 0
	\quad\text{ at }(t_0,x_0),
	\]
	a contradiction. The second case is identical. Hence $\kappa\ge0$, that is,
	\[
	U(t,x) \ge -\frac{\tilde{M}(x^2 + 2dt)}{\ell^2},
	\qquad W(t,x) \ge -\frac{\tilde{M}(x^2 + 2dt)}{\ell^2}
	\qquad (0\le t\le t^*,\ 0\le x\le\ell),
	\]
	and letting $\ell\to\infty$ gives $U\ge0$, $W\ge0$ on
	$[0,t^*]\times[0,\infty)$, i.e. $u\le\overline{u}$ and $v\ge\underline{v}$
	there.

	Finally, $Z:=\overline{u}-u$ satisfies
	$Z_t-d_1Z_{xx}\ge c_{11}Z+c_{12}(\overline{w}-w)\ge c_{11}Z$ on
	$\Omega_{t^*}:=\{0<t\le t^*,\ 0\le x<h(t)\}$ and $Z\ge0$, $Z\not\equiv0$;
	the strong maximum principle gives $Z>0$ in $\Omega_{t^*}$ and, since
	$Z(t^*,h(t^*))=\overline u(t^*,\overline h(t^*))=0$, the Hopf boundary lemma
	gives $Z_x(t^*,h(t^*))<0$, i.e. $\overline u_x(t^*,\overline h(t^*))<u_x(t^*,h(t^*))$.
	Consequently
	\[
	\overline h'(t^*)\ge-\mu\big(A(\overline h(t^*)-ct^*)\big)\overline u_x(t^*,\overline h(t^*))
	>-\mu\big(A(h(t^*)-ct^*)\big)u_x(t^*,h(t^*))=h'(t^*),
	\]
	contradicting \eqref{eq:h}. Therefore $h<\overline{h}$ on $(0,T]$, and the
	above argument applied on $[0,T]$ gives the assertion.

	\textit{Step 2: the case $h_0=\overline{h}(0)$.} For small $\epsilon>0$ let
	$(u_\epsilon,v_\epsilon,h_\epsilon)$ be the solution of \eqref{main} with
	$h_0$ replaced by $(1-\epsilon)h_0$ and $u_0$ replaced by a function
	$u_{0,\epsilon}$ satisfying \eqref{initial} on $[0,(1-\epsilon)h_0]$ with
	$u_{0,\epsilon}\le u_0$ and $u_{0,\epsilon}\to u_0$ in $C^2$. By Step 1,
	$u_\epsilon\le\overline{u}$, $v_\epsilon\ge\underline{v}$ and
	$h_\epsilon<\overline{h}$. By the continuous dependence of the solution of
	\eqref{main} on its data (\cref{th4}), $(u_\epsilon,v_\epsilon,h_\epsilon)\to
	(u,v,h)$ as $\epsilon\to0$, and the conclusion follows by passing to the
	limit.
\end{proof}

\begin{remark}\label{rm:sub_super}
	$(\overline{u},\underline{v},\overline{h})$ and
	$(\underline{u},\overline{v},\underline{h})$ as in \cref{th7} are called an
	upper solution and a lower solution of \eqref{main}, respectively. The same
	proof applies verbatim when $A(\cdot-ct)$ and $\mu(A(\cdot-ct))$ are replaced
	by the constants $a_1$ and $\mu(a_1)$, i.e. \cref{th7} also holds for
	\eqref{main2}; and it applies when the two triples are the solutions of
	\eqref{main} corresponding to two ordered sets of data, which yields the
	monotone dependence of $(u,-v,h)$ on $(u_0,-v_0,h_0)$ and on $\mu$.
\end{remark}

\section{Long-time dynamics: proof of \texorpdfstring{\cref{th:dynamics}}{the dynamics theorem}}\label{sec_3}

By \cref{th6} the free boundary $h$ is strictly increasing, so
\[
h_\infty := \lim_{t\to+\infty}h(t)\in(h_0,+\infty]
\]
is well defined. We treat the two cases $h_\infty<+\infty$ and $h_\infty=+\infty$
separately.

\subsection{The vanishing case}

The following elementary observation is used repeatedly.

\begin{lemma}\label{lem:reduction}
	Suppose that there exists $T_0>0$ such that $h(t)<ct$ for every $t\ge T_0$;
	this is the case, with $T_0:=h_\infty/c$, whenever $h_\infty<+\infty$. Then
	for every $t\ge T_0$,
	\[
	A(x-ct)=a_1 \ \text{ for all } x\in[0,h(t)]\qquad\text{and}\qquad
	\mu(A(h(t)-ct))=\mu(a_1),
	\]
	and consequently $(t,x)\mapsto(u(t+T_0,x),v(t+T_0,x),h(t+T_0))$ is the
	solution of the homogeneous problem \eqref{main2} with $(u_0,v_0,h_0)$
	replaced by $(u(T_0,\cdot),v(T_0,\cdot),h(T_0))$.
\end{lemma}

\begin{proof}
	If $h_\infty<+\infty$ and $t>h_\infty/c$ then $ct>h_\infty\ge h(t)$. In
	either case, $x-ct\le h(t)-ct<0$ for every $x\in[0,h(t)]$, so \eqref{A} gives
	$A(x-ct)=a_1$ and $\mu(A(h(t)-ct))=\mu(a_1)$. Substituting into \eqref{main}
	yields \eqref{main2}.
\end{proof}

\begin{lemma}\label{th1:lm1}
	Let $(u,v,h)$ be the solution of \eqref{main} with initial functions
	satisfying \eqref{initial} and assume \eqref{weak_competition}. If
	$h_\infty<+\infty$, then
	\[
	h_\infty\le R^*,\qquad
	\lim_{t\to\infty} \|u(t,\cdot)\|_{C([0, h(t)])} = 0,\qquad
	\lim_{t\to\infty} v(t,\cdot) = \frac{a_2}{c_2}\ \text{ in } C^2_{\rm loc}([0,\infty)).
	\]
\end{lemma}

\begin{proof}
	By \cref{lem:reduction}, after the time shift $T_0=h_\infty/c$ the triple
	$(u,v,h)$ solves the homogeneous problem \eqref{main2}, whose initial data
	$(u(T_0,\cdot),v(T_0,\cdot),h(T_0))$ satisfy \eqref{initial}: indeed
	$u(T_0,\cdot)>0$ on $[0,h(T_0))$, $u_x(T_0,0)=u(T_0,h(T_0))=0$, $v(T_0,\cdot)>0$,
	$v_x(T_0,0)=0$, and $\liminf_{x\to\infty}v(T_0,x)>0$ by \cref{lem:vlow} below.
	Since for the shifted triple the free boundary still converges to
	$h_\infty<+\infty$, alternative (ii) of \cite[Theorem 1.1]{MR3986328} occurs
	for \eqref{main2}, that is,
	\[
	\lim_{t\to\infty}\|u(t,\cdot)\|_{C([0,h(t)])}=0,\qquad
	\lim_{t\to\infty}v(t,\cdot)=\frac{a_2}{c_2}\ \text{ in }C^2_{\rm loc}([0,\infty)).
	\]
	It remains to prove $h_\infty\le R^*$. Suppose, by contradiction, that
	$h_\infty>R^*$. Since $h$ is increasing with limit $h_\infty$, we may pick
	$T_1>T_0$ with $h(T_1)>R^*$. Applying \cref{lem:reduction} with initial time
	$T_1$, the triple $(u(\cdot+T_1,\cdot),v(\cdot+T_1,\cdot),h(\cdot+T_1))$
	solves \eqref{main2} with initial free boundary $h(T_1)>R^*$. By
	\cite[Theorem 1.2]{MR3986328}, spreading always happens for \eqref{main2} when
	the initial habitat exceeds
	$\frac{\pi}{2}\sqrt{d_1c_2/(a_1c_2-a_2c_1)}=R^*$, whatever the value of
	$\mu(a_1)>0$. Hence $h(t)\to+\infty$, contradicting $h_\infty<+\infty$.
\end{proof}

The following simple lemma, used above and again in \cref{sec_6}, records the
fact that the native population never dies out at infinity.

\begin{lemma}\label{lem:vlow}
	Let $(u,v,h)$ be the solution of \eqref{main} with \eqref{initial}. Let
	$(\hat u,\hat v)$ solve the kinetic system
	\begin{equation}\label{ode}
		\hat u' = (a_1 - b_1\hat u - c_1 \hat v )\hat u,\quad
		\hat v' = (a_2 -b_2 \hat u - c_2 \hat v)\hat v,\quad
		(\hat{u}(0), \hat{v}(0)) = \big(\|u_0\|_\infty, \sigma_0^{\max}\big),
	\end{equation}
	with $\sigma_0^{\max}$ as in \eqref{v0inf}. Then
	\[
	u(t,x)\le\hat u(t),\qquad v(t,x)\ge\hat v(t)\qquad (t>0,\ x\ge0),
	\]
	and $(\hat u(t),\hat v(t))\to(u^*,v^*)$ as $t\to\infty$. In particular, for
	every $\varepsilon>0$ there is $T_\varepsilon>0$ with
	\begin{equation}\label{eq:uv_global}
		u(t,x)\le u^*+\varepsilon,\qquad v(t,x)\ge v^*-\varepsilon
		\qquad (t\ge T_\varepsilon,\ x\ge0).
	\end{equation}
\end{lemma}

\begin{proof}
	Since $A(\cdot)\le a_1$, the triple $(\hat u,\hat v,+\infty)$ is an upper
	solution in the sense of \cref{th7}(a) (with the first alternative at $x=0$,
	both derivatives being $0$), because $\hat u,\hat v$ do not depend on $x$ and
	$\hat u(0)=\|u_0\|_\infty\ge u_0$, $\hat v(0)=\sigma_0^{\max}\le v_0$. The
	comparison principle gives the stated inequalities. The convergence
	$(\hat u,\hat v)\to(u^*,v^*)$ is the classical global stability of the
	coexistence state of a weakly competing Lotka--Volterra system under
	\eqref{weak_competition}; see e.g. \cite[Chapter 4]{MR1319817}.
\end{proof}

\subsection{The spreading case}

\begin{lemma}\label{th1:lm2}
	Let $(u,v,h)$ be the solution of \eqref{main} with \eqref{initial} and assume
	\eqref{weak_competition}. If $h_\infty=+\infty$, then
	\[
	\lim_{t\to\infty} (u(t,\cdot), v(t,\cdot)) = (u^*, v^*)\ \text{ in } C^2_{\rm loc}([0, \infty)).
	\]
\end{lemma}

\begin{proof}
	\textit{Step 0: upper bound for $u$ and lower bound for $v$.} Since
	$A(\cdot)\le a_1$ and $\mu(A(\cdot))\le\mu(a_1)$ while $-u_x(t,h(t))>0$, the
	triple $(u,v,h)$ is a lower solution of \eqref{main2} in the sense of
	\cref{th7}(b) (first alternative at $x=0$). Let
	$(\overline{u},\underline{v},\overline{h})$ denote the solution of
	\eqref{main2} with the same initial data. Then
	\begin{equation}\label{eq:comp_main2}
		u\le\overline{u},\qquad v\ge\underline{v},\qquad h\le\overline{h}.
	\end{equation}
	Since $h_\infty=+\infty$ we get $\overline{h}(t)\to+\infty$, so alternative
	(i) of \cite[Theorem 1.1]{MR3986328} holds for \eqref{main2}, that is,
	$(\overline{u},\underline{v})\to(u^*,v^*)$ in $C^2_{\rm loc}([0,\infty))$.
	Combining this with \eqref{eq:comp_main2},
	\begin{equation}\label{uv1}
		\limsup_{t \to +\infty} u(t, x) \le u^*, \qquad \liminf_{t \to +\infty} v(t, x) \ge v^*
	\end{equation}
	uniformly on compact subsets of $[0,\infty)$. It therefore suffices to prove
	\begin{equation}\label{uv2}
		\liminf_{t \to +\infty} u(t, x) \ge u^*, \qquad \limsup_{t \to +\infty} v(t, x) \le v^*
	\end{equation}
	locally uniformly in $[0,\infty)$.

	\textit{Step 1: an iteration scheme.} Set
	\[
	\overline{v}_0:=\frac{a_2}{c_2},\qquad
	\underline{u}_{n+1}:=\frac{a_1-c_1\overline{v}_n}{b_1},\qquad
	\overline{v}_{n+1}:=\frac{a_2-b_2\underline{u}_{n+1}}{c_2}\qquad(n\ge0).
	\]
	The map $\underline u\mapsto \frac{a_1}{b_1}-\frac{c_1}{b_1}
	\big(\frac{a_2}{c_2}-\frac{b_2}{c_2}\underline u\big)
	=\frac{a_1c_2-a_2c_1}{b_1c_2}+\frac{b_2c_1}{b_1c_2}\underline u$ is affine with
	slope $\frac{b_2c_1}{b_1c_2}\in(0,1)$, because \eqref{weak_competition} gives
	$c_1/c_2<b_1/b_2$; its unique fixed point is
	$\frac{a_1c_2-a_2c_1}{b_1c_2-b_2c_1}=u^*$. Since
	$\underline{u}_1=\frac{a_1c_2-a_2c_1}{b_1c_2}\in(0,u^*)$ by
	\eqref{ineq_star}, it follows that $\{\underline{u}_n\}$ is strictly
	increasing with $\underline{u}_n\uparrow u^*$ and $\{\overline{v}_n\}$ is
	strictly decreasing with $\overline{v}_n\downarrow v^*$; in particular
	\begin{equation}\label{eq:pos_iter}
		0<\underline{u}_n<u^*\ \ (n\ge1),
		\qquad
		v^*<\overline{v}_n\le\frac{a_2}{c_2}
		\ \text{ and }\
		a_1-c_1\overline{v}_n\ge a_1-\frac{a_2c_1}{c_2}>0\ \ (n\ge0),
	\end{equation}
	the last inequality by \eqref{ineq_star}.
	We claim that for every $n\ge0$,
	\begin{equation}\label{induc}
		\liminf_{t \to +\infty} u(t,x) \ge \underline{u}_n,  \qquad \limsup_{t \to +\infty} v(t,x) \le \overline{v}_n
	\end{equation}
	uniformly on compact subsets of $[0,\infty)$, with the convention
	$\underline u_0:=0$. Granting \eqref{induc}, letting $n\to\infty$ gives
	\eqref{uv2} and hence, together with \eqref{uv1}, the locally uniform
	convergence $(u,v)\to(u^*,v^*)$. Interior parabolic estimates (applied on
	$[T,T+2]\times[0,\ell]$ for $\ell$ fixed and $T\to\infty$, using the uniform
	bounds of \cref{th6}) then upgrade this convergence to $C^2_{\rm loc}([0,\infty))$,
	which is the assertion of the lemma.

	\textit{Step 2: \eqref{induc} for $n=0$.} The inequality
	$\liminf u\ge\underline u_0=0$ is trivial. For $v$, the comparison principle
	applied to $v_t-d_2v_{xx}\le v(a_2-c_2v)$ and to the solution $V_1$ of
	$V_1'=V_1(a_2-c_2V_1)$, $V_1(0)=\|v_0\|_\infty$, gives $v(t,x)\le V_1(t)$ for
	all $x\ge0$, $t>0$, whence
	\begin{equation}\label{v1}
		\limsup_{t \to +\infty}{v(t,x)} \le \frac{a_2}{c_2} = \overline{v}_0
		\quad\text{ uniformly for } x \in [0,\infty).
	\end{equation}

	\textit{Step 3: from $n$ to $n+1$.} Assume \eqref{induc} holds for some
	$n\ge0$. Fix $\varepsilon>0$ so small that
	\begin{equation}\label{eq:eps_small}
		a_1-c_1(\overline v_n+\varepsilon)>0
		\qquad\text{and}\qquad
		a_2-b_2(\underline u_{n+1}-\varepsilon)>0,
	\end{equation}
	which is possible by \eqref{eq:pos_iter}, and fix
	\begin{equation}\label{eq:ell_large}
		\ell>\max\left\{h_0,\ \frac{\pi}{2}\sqrt{\frac{d_1}{a_1-c_1(\overline v_n+\varepsilon)}},\
		\frac{\pi}{2}\sqrt{\frac{d_2}{a_2-b_2(\underline u_{n+1}-\varepsilon)}}\right\}.
	\end{equation}
	Since $h_\infty=+\infty$ and by the induction hypothesis, there exists
	$t_n>0$ such that
	\begin{equation}\label{eq:tn}
		h(t)>\ell,\quad ct>\ell,\quad
		u(t,x)\ge\underline{u}_n-\varepsilon,\quad
		v(t,x)\le \overline{v}_n+\varepsilon
		\qquad (t>t_n,\ x\in[0,\ell]).
	\end{equation}
	Because $ct>\ell$ we have $A(x-ct)=a_1$ for all $x\in[0,\ell]$ and $t>t_n$, so
	$u$ satisfies
	\[
	\begin{cases}
		u_t - d_1u_{xx} \ge \big(a_1 - b_1 u - c_1 (\overline{v}_n+\varepsilon)\big) u , & 0 < x< \ell, ~ t > t_n,\\
		u_x(t, 0) = 0, \quad u(t, \ell) > 0, & t > t_n .
	\end{cases}
	\]
	Hence $u$ is an upper solution of
	\[
	\begin{cases}
		\hat{u}_t = d_1 \hat{u}_{xx} + \big(a_1 - b_1 \hat{u} - c_1 (\overline{v}_n + \varepsilon)\big)\hat{u}, & 0 < x < \ell, ~ t > t_n,\\
		\hat{u}_x(t,0) = 0, \quad \hat{u}(t,\ell) = 0, & t > t_n,\\
		\hat{u}(t_n,x) = u(t_n,x), & 0 \le x \le \ell,
	\end{cases}
	\]
	so that $u(t,x)\ge\hat u(t,x)$ on $[0,\ell]$ for $t>t_n$. By the choice
	\eqref{eq:ell_large} of $\ell$, the principal eigenvalue of
	$-d_1\partial_{xx}$ on $(0,\ell)$ with Neumann condition at $0$ and Dirichlet
	condition at $\ell$, namely $d_1\pi^2/(4\ell^2)$, is smaller than
	$a_1-c_1(\overline v_n+\varepsilon)$; hence the above logistic problem has a
	unique positive steady state $\tilde u_\ell$, $\hat u(t,\cdot)\to\tilde u_\ell$
	uniformly on $[0,\ell]$ as $t\to\infty$, and
	$\tilde{u}_\ell\to\frac{a_1-c_1(\overline v_n+\varepsilon)}{b_1}$ locally
	uniformly on $[0,\infty)$ as $\ell\to\infty$ (see e.g.
	\cite[Lemma 2.2]{DU_MA_2001}). Since $\ell$ may be taken arbitrarily large and
	$\varepsilon>0$ arbitrarily small, we conclude
	\begin{equation}\label{eq:un1}
		\liminf_{t \to +\infty}{u(t,x)} \ge \frac{a_1-c_1\overline{v}_n}{b_1} = \underline{u}_{n+1}
		\quad\text{ locally uniformly in }[0,\infty).
	\end{equation}
	The bound for $v$ is obtained in exactly the same way. Indeed, by
	\eqref{eq:un1} there is $t_n'>t_n$ with
	$u(t,x)\ge\underline{u}_{n+1}-\varepsilon$ for $x\in[0,\ell]$, $t>t_n'$, and
	therefore
	\[
	\begin{cases}
		v_t - d_2v_{xx} \le \big(a_2 - b_2(\underline{u}_{n+1}-\varepsilon) - c_2 v\big) v , & 0 < x< \ell, ~ t > t_n',\\
		v_x(t, 0) = 0, \quad v(t, \ell) \le M_2, & t > t_n' ,
	\end{cases}
	\]
	so that $v\le\hat v$ on $[0,\ell]\times(t_n',\infty)$, where $\hat v$ solves
	the corresponding logistic problem with boundary value $M_2$ at $x=\ell$ and
	initial value $v(t_n',\cdot)$. Again by \eqref{eq:ell_large} this problem has a
	unique positive steady state, to which $\hat v(t,\cdot)$ converges, and that
	steady state tends to $\frac{a_2-b_2(\underline u_{n+1}-\varepsilon)}{c_2}$
	locally uniformly as $\ell\to\infty$. Letting $\ell\to\infty$ and
	$\varepsilon\to0$,
	\[
	\limsup_{t \to +\infty}{v(t,x)} \le \frac{a_2-b_2\underline{u}_{n+1}}{c_2} = \overline{v}_{n+1}
	\quad\text{ locally uniformly in }[0,\infty),
	\]
	which completes the induction and the proof.
\end{proof}

\begin{proof}[Proof of \cref{th:dynamics}]
	By \cref{th6}, $h_\infty$ exists in $(h_0,+\infty]$. If $h_\infty<+\infty$,
	\cref{th1:lm1} gives alternative (ii); if $h_\infty=+\infty$,
	\cref{th1:lm2} gives alternative (i). The two alternatives are mutually
	exclusive.
\end{proof}

\section{Criteria for spreading and vanishing: proof of \texorpdfstring{\cref{th:criteria}}{the criteria theorem}}\label{sec_4}

\begin{proof}[Proof of \cref{th:criteria}(i)]
	Assume $h_0\ge R^*$ and suppose that vanishing occurs. By \cref{th1:lm1} we
	would have $h_\infty\le R^*\le h_0$, contradicting $h_\infty>h_0$. Hence
	spreading occurs.
\end{proof}

\begin{proof}[Proof of \cref{th:criteria}(ii)]
	Assume $h_0<R_0$ and write $u_0=\sigma\phi$.

	\textit{Monotonicity in $\sigma$.} If $0<\sigma_1<\sigma_2$ then
	$\sigma_1\phi\le\sigma_2\phi$, and by \cref{rm:sub_super} the triple
	$(u^{\sigma_1},v^{\sigma_1},h^{\sigma_1})$ is a lower solution of the problem
	solved by $(u^{\sigma_2},v^{\sigma_2},h^{\sigma_2})$; the comparison principle
	yields $h^{\sigma_1}(t)\le h^{\sigma_2}(t)$ for all $t\ge0$. Hence, if
	spreading occurs for $\sigma_1$, then it occurs for every $\sigma\ge\sigma_1$.
	Consequently the set
	$\Sigma:=\{\sigma>0:\ \text{vanishing occurs for } u_0=\sigma\phi\}$ is an
	interval of the form $(0,\sigma_0]$ or $(0,\sigma_0)$ with
	$\sigma_0\in[0,+\infty]$, and it remains to rule out $\sigma_0=0$ and to check
	that the endpoint belongs to $\Sigma$.

	\textit{$\sigma_0>0$.} Since $A(\cdot)\le a_1$ and $v>0$, the pair $(u,h)$
	satisfies
	\[
	u_t-d_1u_{xx}\le u(a_1-b_1u)\ \ \text{ in }\{0<x<h(t)\},\qquad
	h'(t)\le -\mu(a_1)u_x(t,h(t)),
	\]
	together with $u_x(t,0)=0$ and $u(t,h(t))=0$. Thus $(u,h)$ is a lower solution
	of the single species free boundary problem
	\begin{equation}\label{DL}
		\begin{cases}
			U_t = d_1U_{xx}+U(a_1-b_1U), & t>0,\ 0<x<H(t),\\
			U_x(t,0)=U(t,H(t))=0,\ \ H'(t)=-\mu(a_1)U_x(t,H(t)), & t>0,\\
			H(0)=h_0,\ U(0,x)=\sigma\phi(x), & 0\le x\le h_0,
		\end{cases}
	\end{equation}
	so $h(t)\le H(t)$ for all $t$. Assume
	$h_0<R_0=\frac{\pi}{2}\sqrt{d_1/a_1}$. The upper solution constructed in the
	proof of \cite[Lemma 3.8]{MR2607347} then applies to \eqref{DL}: it forces
	$H_\infty<+\infty$ as soon as $\mu(a_1)M$ is small enough, where $M$ is any
	constant with $\sigma\phi(x)\le M\cos\big(\frac{\pi x}{2h_0(1+\delta/2)}\big)$
	on $[0,h_0]$ and $\delta$ depends only on $d_1,a_1,h_0$. Since $M$ may be
	taken proportional to $\sigma\|\phi\|_{\infty}$, vanishing occurs for
	\eqref{DL} whenever $\sigma\|\phi\|_{\infty}$ is small enough, $\mu(a_1)$
	being fixed. In that case $H_\infty<+\infty$, so
	$h_\infty\le H_\infty<+\infty$ and vanishing occurs for \eqref{main} as well,
	by \cite[Lemma 3.1]{MR2607347}. Therefore $\sigma_0>0$.

	\textit{The endpoint.} Let $\sigma_0<+\infty$ and suppose spreading occurred
	for $\sigma=\sigma_0$. Then $h^{\sigma_0}(T)>R^*$ for some $T>0$. The solution
	of \eqref{main} depends continuously on its initial data on the fixed time
	interval $[0,T]$: this follows from \cref{th4} and a standard continuation
	argument, the solution being global with bounds that are uniform for
	$\sigma$ in a compact set by \cref{th6}. Hence there is
	$\sigma<\sigma_0$ with $h^{\sigma}(T)>R^*$, and
	\cref{th1:lm1} forces spreading for that $\sigma$, contradicting
	$\sigma\in\Sigma$. Hence $\sigma_0\in\Sigma$, which is the assertion.
\end{proof}

\section{Semi-wave solutions}\label{sec_5}

Throughout this section we assume \eqref{weak_competition}. We first recall the
two results of \cite{MR3986328} that we shall use.

\begin{theorem}[see {\cite[Proposition 1.4]{MR3986328}} and \cref{rm:convert}]\label{th:c*}
	There exists a critical speed
	\[
	c^*\ \ge\ 2\sqrt{d_1\left(a_1-\frac{a_2c_1}{c_2}\right)}\ >0
	\]
	such that the problem
	\begin{equation}\label{tw}
		\begin{cases}
			-c \Psi' - d_1 \Psi'' = (a_1 - b_1\Psi - c_1\Phi)\Psi, & -\infty < x < +\infty,\\
			-c \Phi' - d_2 \Phi'' = (a_2 - b_2\Psi - c_2\Phi)\Phi, & -\infty < x < +\infty,\\
			\Psi(-\infty) = u^*, ~ \Psi(+\infty)=0, ~ \Psi'(x)<0 \text{ for } x \in \mathbb{R},\\
			\Phi(-\infty) = v^*, ~ \Phi(+\infty)=\frac{a_2}{c_2}, ~ \Phi'(x)>0 \text{ for } x\in \mathbb{R}
		\end{cases}
	\end{equation}
	has a solution when $c\ge c^*$ and has no solution when $c<c^*$.
\end{theorem}

\begin{theorem}[see {\cite[Theorem 1.5]{MR3986328}} and \cref{rm:convert}]\label{th:semi-wave0}
	For each $c\in[0,c^*)$ the problem
	\begin{equation}\label{semiwave0}
		\begin{cases}
			-c \psi' - d_1 \psi'' = (a_1 - b_1\psi - c_1\varphi)\psi, & -\infty < x < 0,\\
			-c \varphi' - d_2 \varphi'' = (a_2 - b_2\psi - c_2\varphi)\varphi, & -\infty < x < +\infty,\\
			\psi(-\infty) = u^*, ~ \psi(x)=0 \text{ for } x\ge 0, ~ \psi'(x)<0 \text{ for } x\le 0,\\
			\varphi(-\infty) = v^*, ~ \varphi(+\infty)=\frac{a_2}{c_2}, ~ \varphi'(x)>0 \text{ for } x\in \mathbb{R}
		\end{cases}
	\end{equation}
	has a unique solution $(\Psi_c, \Phi_c) \in [C(\mathbb{R})\cap C^2((-\infty, 0])] \times C^2(\mathbb{R})$,
	and it has no such solution for $c\ge c^*$. The map
	$c\mapsto(\Psi_c,\Phi_c)$ is continuous from $[0,c^*)$ into
	$C^2_{\rm loc}((-\infty,0])\times C^2_{\rm loc}(\mathbb{R})$; in particular
	$c\mapsto\Psi_c'(0)$ is continuous on $[0,c^*)$. Moreover there is a unique
	$c_0\in(0,c^*)$ with
	\begin{equation}\label{c0}
		-\mu(a_1) \Psi_{c_0}'(0) = c_0 ,
	\end{equation}
	and $c_0$ depends continuously on $a_1,a_2,b_1,b_2,c_1,c_2,d_1,d_2$ and
	$\mu(a_1)$.
\end{theorem}

\begin{remark}\label{rm:convert}
	Theorems~\ref{th:c*} and \ref{th:semi-wave0} are stated in \cite{MR3986328}
	for the non-dimensional system obtained from \eqref{main2} by the scaling
	$\hat u(t,x)=\frac{b_1}{a_1}u\big(t/a_1,\sqrt{d_1/a_1}\,x\big)$,
	$\hat v(t,x)=\frac{c_2}{a_2}v\big(t/a_1,\sqrt{d_1/a_1}\,x\big)$, in which the
	competition coefficients become $k=\frac{a_2c_1}{a_1c_2}$ and
	$h=\frac{a_1b_2}{a_2b_1}$, so that \eqref{weak_competition} is equivalent to
	$k,h\in(0,1)$. Speeds are multiplied by $\sqrt{a_1d_1}$ under this scaling,
	so the bound $c^*\ge2\sqrt{1-k}$ of \cite[Proposition 1.4]{MR3986328} becomes
	$c^*\ge2\sqrt{a_1d_1(1-k)}=2\sqrt{d_1(a_1-a_2c_1/c_2)}$, as stated in
	\cref{th:c*}. The continuity of $c\mapsto(\Psi_c,\Phi_c)$ follows from the
	uniqueness statement together with standard elliptic estimates and a
	compactness argument, exactly as in \cite[Theorem 1.3(ii)]{MR3609207}.
\end{remark}

The main result of this section supplies the semi-waves adapted to the shifting
environment. It is the analogue of \cite[Proposition 1.1]{HU20205931} for the
competition system, and of the corresponding result of \cite{DLN2026a} in the
weak--strong regime.

\begin{theorem}\label{th:semi-wave}
	Assume \eqref{weak_competition} and $0<c\le c_0$. Then the following hold.
	\begin{itemize}
		\item[(i)] For every $L\ge0$ the problem
		\begin{equation}\label{eqlm1}
			\begin{cases}
				-c \psi' - d_1 \psi'' = (A(x) - b_1\psi - c_1\varphi)\psi, & -\infty < x < L,\\
				-c \varphi' - d_2 \varphi'' = (a_2 - b_2\psi - c_2\varphi)\varphi, & -\infty < x < +\infty,\\
				\psi(-\infty) = u^*, ~ \psi(x)=0 \text{ for } x\ge L, ~ \psi'(x)<0 \text{ for } x\le L,\\
				\varphi(-\infty) = v^*, ~ \varphi(+\infty)=\frac{a_2}{c_2}, ~ \varphi'(x)>0 \text{ for } x\in \mathbb{R}
			\end{cases}
		\end{equation}
		has a unique solution
		$(\psi_L,\varphi_L) \in [C(\mathbb{R})\cap C^2((-\infty, L])] \times C^2(\mathbb{R})$.
		\item[(ii)] The map $L \mapsto \psi'_L(L)$ is strictly increasing on
		$[0, +\infty)$ and $\lim_{L\to+\infty}\psi'_L(L)=0$.
		\item[(iii)] There exists a unique $L_0 \geq 0$ such that
		$-\mu(A(L_0))\psi'_{L_0}(L_0) = c$. Moreover $L_0 = 0$ if and only if
		$c = c_0$, in which case $(\psi_0,\varphi_0) \equiv (\Psi_{c_0}, \Phi_{c_0})$.
	\end{itemize}
\end{theorem}

The proof occupies the rest of this section. It is convenient to transform
\eqref{eqlm1} into a cooperative system. Setting
\begin{equation}\label{cooper}
	\Theta_1:=\psi,\qquad \Theta_2:=\frac{a_2}{c_2}-\varphi ,
\end{equation}
a direct computation shows that \eqref{eqlm1} is equivalent to
\begin{equation}\label{eqcoop}
	\begin{cases}
		d_1\Theta_1''+c\Theta_1'+F_1(x,\Theta)=0, & -\infty<x<L,\\
		d_2\Theta_2''+c\Theta_2'+F_2(\Theta)=0, & -\infty<x<+\infty,
	\end{cases}
\end{equation}
where
\[
F_1(x,\Theta):=\Big(A(x)-\frac{a_2c_1}{c_2}-b_1\Theta_1+c_1\Theta_2\Big)\Theta_1,
\qquad
F_2(\Theta):=\big(b_2\Theta_1-c_2\Theta_2\big)\Big(\frac{a_2}{c_2}-\Theta_2\Big).
\]
On the rectangle $\mathcal R:=[0,u^*]\times[0,\frac{a_2}{c_2}-v^*]$ we have
\[
\partial_{\Theta_2}F_1=c_1\Theta_1\ge0,
\qquad
\partial_{\Theta_1}F_2=b_2\Big(\frac{a_2}{c_2}-\Theta_2\Big)\ge0 ,
\]
so \eqref{eqcoop} is a cooperative system on $\mathcal R$, and $F_1,F_2$ are
locally Lipschitz. Note that $(u^*,\frac{a_2}{c_2}-v^*)$ is the unique zero of
$(F_1(\cdot),F_2)$ in $\mathcal R$ with both components positive when $A\equiv a_1$.

\subsection{Existence}

\begin{lemma}\label{lm:supersub}
	Let $0<c\le c_0$ and $L\ge0$. Then
	\[
	\overline\Theta:=\Big(u^*,\ \frac{a_2}{c_2}-v^*\Big),
	\qquad
	\underline\Theta:=\Big(\Psi_{c_0},\ \frac{a_2}{c_2}-\Phi_{c_0}\Big)
	\]
	are, respectively, an upper and a lower solution of \eqref{eqcoop} on
	$(-\infty,L)\times \mathbb{R}$, and $\underline\Theta\le\overline\Theta$ componentwise.
\end{lemma}

\begin{proof}
	Since $\Psi_{c_0}'<0$ on $(-\infty,0]$, $\Psi_{c_0}(-\infty)=u^*$ and
	$\Psi_{c_0}\equiv0$ on $[0,\infty)$, we have $0\le\Psi_{c_0}<u^*$; since
	$\Phi_{c_0}'>0$, $\Phi_{c_0}(-\infty)=v^*$ and $\Phi_{c_0}(+\infty)=a_2/c_2$,
	we have $v^*<\Phi_{c_0}<a_2/c_2$. Hence
	$\underline\Theta\le\overline\Theta$ and both take values in $\mathcal R$.

	For $\overline\Theta$ we compute, using $A(x)\le a_1$,
	$a_1-b_1u^*-c_1v^*=0$ and $a_2-b_2u^*-c_2v^*=0$,
	\[
	F_1\Big(x,\ u^*,\ \frac{a_2}{c_2}-v^*\Big)=\big(A(x)-b_1u^*-c_1v^*\big)u^*
	=\big(A(x)-a_1\big)u^*\le0,
	\qquad F_2\Big(u^*,\frac{a_2}{c_2}-v^*\Big)=0 ,
	\]
	so that $d_i\overline\Theta_i''+c\overline\Theta_i'+F_i\le0$: $\overline\Theta$
	is an upper solution.

	For $\underline\Theta$, recall that $(\Psi_{c_0},\Phi_{c_0})$ solves
	\eqref{semiwave0} with $c$ replaced by $c_0$. Hence, for $x<0$,
	\[
	d_1\Psi_{c_0}''+c\Psi_{c_0}'
	= d_1\Psi_{c_0}''+c_0\Psi_{c_0}'+(c-c_0)\Psi_{c_0}'
	= -\big(a_1-b_1\Psi_{c_0}-c_1\Phi_{c_0}\big)\Psi_{c_0}+(c-c_0)\Psi_{c_0}' ,
	\]
	and since $c\le c_0$ and $\Psi_{c_0}'<0$ we get $(c-c_0)\Psi_{c_0}'\ge0$;
	using also $A(x)=a_1$ for $x\le0$ this gives
	$d_1\underline\Theta_1''+c\underline\Theta_1'+F_1(x,\underline\Theta)\ge0$
	on $(-\infty,0)$. On $(0,L)$ we have $\underline\Theta_1\equiv0$ and both
	sides vanish; the function $\underline\Theta_1=\Psi_{c_0}$ has a corner at
	$x=0$ with $\Psi_{c_0}'(0^-)<0=\Psi_{c_0}'(0^+)$, so $\underline\Theta_1$ is a
	weak lower solution in the sense of \cite[Definition 2.4]{MR3986328}, which is
	admissible for the monotone iteration scheme. Similarly, for $x\in \mathbb{R}$,
	\[
	d_2\Big(\frac{a_2}{c_2}-\Phi_{c_0}\Big)''+c\Big(\frac{a_2}{c_2}-\Phi_{c_0}\Big)'
	= -\big[d_2\Phi_{c_0}''+c_0\Phi_{c_0}'\big]-(c-c_0)\Phi_{c_0}'
	= \big(a_2-b_2\Psi_{c_0}-c_2\Phi_{c_0}\big)\Phi_{c_0}-(c-c_0)\Phi_{c_0}' ,
	\]
	and $-(c-c_0)\Phi_{c_0}'\ge0$ because $c\le c_0$ and $\Phi_{c_0}'>0$; since
	$\big(a_2-b_2\Psi_{c_0}-c_2\Phi_{c_0}\big)\Phi_{c_0}
	=-F_2(\underline\Theta)$, we obtain
	$d_2\underline\Theta_2''+c\underline\Theta_2'+F_2(\underline\Theta)\ge0$.
\end{proof}

\begin{proof}[Proof of \cref{th:semi-wave}(i), existence]
	Fix $L\ge0$. For $n>L$ consider the truncated problem
	\[
	\begin{cases}
		d_1\Theta_1''+c\Theta_1'+F_1(x,\Theta)=0, & -n<x<L,\\
		d_2\Theta_2''+c\Theta_2'+F_2(\Theta)=0, & -n<x<n,\\
		\Theta_1(-n)=\overline\Theta_1(-n),\quad \Theta_1(x)=0\ \ (L\le x\le n),\\
		\Theta_2(-n)=\overline\Theta_2(-n),\quad \Theta_2(n)=\overline\Theta_2(n).
	\end{cases}
	\]
	Since \eqref{eqcoop} is cooperative on $\mathcal R$ and, by
	\cref{lm:supersub}, $\underline\Theta\le\overline\Theta$ are ordered weak
	lower and upper solutions, the standard monotone iteration scheme for
	quasimonotone elliptic systems converges to a solution $\Theta^n$ with
	$\underline\Theta\le\Theta^n\le\overline\Theta$ on the truncated domain.
	Schauder estimates give bounds on $\Theta^n$ in $C^{2,\alpha}$ over compact
	subsets that are independent of $n$, so a diagonal extraction yields a
	classical solution $\Theta=(\Theta_1,\Theta_2)$ of \eqref{eqcoop} on
	$(-\infty,L]\times \mathbb{R}$ with
	\begin{equation}\label{s0}
		\Psi_{c_0}(x) \le \psi_L(x) \le u^*,
		\qquad
		v^* \le \varphi_L(x) \le \Phi_{c_0}(x) < \frac{a_2}{c_2}
		\qquad (x\in \mathbb{R}),
	\end{equation}
	where $(\psi_L,\varphi_L)$ is recovered from $\Theta$ through
	\eqref{cooper}. In particular
	\begin{equation}\label{s1}
		\psi_L(-\infty) = u^*, \qquad \varphi_L(-\infty) = v^* .
	\end{equation}
	Moreover $\psi_L>0$ on $(-\infty,L)$: this holds on $(-\infty,0)$ because
	$\psi_L\ge\Psi_{c_0}>0$ there, and on $[0,L)$ because otherwise the strong
	maximum principle applied to the $\psi$-equation would force $\psi_L\equiv0$
	near an interior zero and hence on $(-\infty,L)$ by unique continuation,
	contradicting $\psi_L\ge\Psi_{c_0}>0$ on $(-\infty,0)$. By the Hopf boundary
	lemma, $\psi_L'(L)<0$. Similarly, $\varphi_L>v^*$ on $\mathbb{R}$ by the strong
	maximum principle.

	It remains to prove the monotonicity statements in \eqref{eqlm1} and
	$\varphi_L(+\infty)=a_2/c_2$; this is done in
	Lemmas~\ref{lm:phi_infty}--\ref{lm:mono} below.
\end{proof}

\begin{lemma}\label{lm:phi_infty}
	$\varphi_L'>0$ on $(L,+\infty)$ and $\varphi_L(+\infty)=\frac{a_2}{c_2}$.
\end{lemma}

\begin{proof}
	For $x>L$ we have $\psi_L(x)=0$, hence
	\begin{equation}\label{s3}
		-c \varphi_L' - d_2 \varphi_L'' = (a_2 - c_2\varphi_L)\varphi_L > 0
		\qquad (x>L),
	\end{equation}
	because $v^*\le\varphi_L<a_2/c_2$ by \eqref{s0}. Equivalently
	$d_2\varphi_L''+c\varphi_L'<0$, and multiplying by
	$\frac{1}{d_2}e^{\frac{c}{d_2}x}$ gives
	\[
	\Big(e^{\frac{c}{d_2}x}\varphi_L'\Big)'
	=\frac{1}{d_2}e^{\frac{c}{d_2}x}\big(d_2\varphi_L''+c\varphi_L'\big)<0
	\qquad (x>L),
	\]
	so $x\mapsto e^{\frac{c}{d_2}x}\varphi_L'(x)$ is strictly decreasing on
	$(L,+\infty)$.

	Suppose $\varphi_L'(x_1)\le0$ for some $x_1>L$. Then
	$e^{\frac{c}{d_2}x}\varphi_L'(x)<e^{\frac{c}{d_2}x_1}\varphi_L'(x_1)\le0$
	for every $x>x_1$, so $\varphi_L'<0$ on $(x_1,\infty)$ and $\varphi_L$ is
	strictly decreasing there; being bounded below by $v^*>0$, it has a limit
	$m\in[v^*,\varphi_L(x_1))\subset(0,a_2/c_2)$. By \eqref{s0} and elliptic
	estimates, $\varphi_L$ is bounded in $C^{2,\alpha}$ on $(L+1,\infty)$, whence
	$\varphi_L'(x)\to0$ and $\varphi_L''(x)\to0$ along a sequence $x\to\infty$.
	Passing to the limit in \eqref{s3} gives $(a_2-c_2m)m=0$, which is impossible
	since $0<m<a_2/c_2$. Therefore $\varphi_L'>0$ on $(L,+\infty)$, $\varphi_L$ is
	increasing and bounded there, and the same limiting argument applied to its
	limit $m\in(0,a_2/c_2]$ gives $(a_2-c_2m)m=0$, i.e. $m=\frac{a_2}{c_2}$.
\end{proof}

\begin{lemma}\label{lm:asym}
	Let $(\psi,\varphi)$ be any solution of \eqref{eqlm1}, and write
	$\psi_L,\varphi_L$ for it. Let $\lambda_1$ be the smallest positive root of
	\begin{equation}\label{charac}
		\Lambda(\lambda):=\big(d_1\lambda^2+c\lambda-b_1u^*\big)\big(d_2\lambda^2+c\lambda-c_2v^*\big)-b_2c_1u^*v^* .
	\end{equation}
	Then $\Lambda$ has exactly two positive roots $0<\lambda_1<\lambda_2$, and
	there are constants $k_1,k_2>0$ (depending on $L$) such that, as
	$x\to-\infty$,
	\begin{equation}\label{asym}
		u^*-\psi_L(x)=\big(k_1+o(1)\big)e^{\lambda_1x},
		\qquad
		\varphi_L(x)-v^*=\big(k_2+o(1)\big)e^{\lambda_1x}.
	\end{equation}
	In particular $\lambda_1$ does not depend on $L$.
\end{lemma}

\begin{proof}
	Write $p(\lambda):=d_1\lambda^2+c\lambda-b_1u^*$ and
	$q(\lambda):=d_2\lambda^2+c\lambda-c_2v^*$, and let
	$\alpha_-<0<\alpha_+$ be the roots of $p$. Then
	$\Lambda(0)=b_1c_2u^*v^*-b_2c_1u^*v^*=u^*v^*(b_1c_2-b_2c_1)>0$ by
	\eqref{weak_competition}, $\Lambda(\alpha_\pm)=-b_2c_1u^*v^*<0$ and
	$\Lambda(\pm\infty)=+\infty$. Hence $\Lambda$ has one root in
	$(0,\alpha_+)$, one in $(\alpha_+,+\infty)$, one in $(\alpha_-,0)$ and one in
	$(-\infty,\alpha_-)$: exactly two positive roots
	$0<\lambda_1<\alpha_+<\lambda_2$.

	Set $\hat\psi:=u^*-\psi_L$ and $\hat\varphi:=\varphi_L-v^*$. They tend
	to $0$ as $x\to-\infty$ because $\psi_L(-\infty)=u^*$ and
	$\varphi_L(-\infty)=v^*$ in \eqref{eqlm1}, and they are positive because
	\eqref{eqlm1} also requires $\psi_L'<0$ on $(-\infty,L]$ and
	$\varphi_L'>0$ on $\mathbb{R}$. \textup{(}For the solution produced in
	part \textup{(i)} the positivity is also immediate from \eqref{s0}.\textup{)} Since
	$A\equiv a_1$ on $(-\infty,0]$ and
	$a_1-b_1\psi_L-c_1\varphi_L=b_1\hat\psi-c_1\hat\varphi$,
	$a_2-b_2\psi_L-c_2\varphi_L=b_2\hat\psi-c_2\hat\varphi$, the pair
	$(\hat\psi,\hat\varphi)$ satisfies, on $(-\infty,0]$,
	\begin{equation}\label{lin}
		\begin{cases}
			d_1\hat\psi''+c\hat\psi'-b_1u^*\hat\psi+c_1u^*\hat\varphi
			=-\big(b_1\hat\psi-c_1\hat\varphi\big)\hat\psi ,\\[2pt]
			d_2\hat\varphi''+c\hat\varphi'+b_2v^*\hat\psi-c_2v^*\hat\varphi
			=-\big(b_2\hat\psi-c_2\hat\varphi\big)\hat\varphi ,
		\end{cases}
	\end{equation}
	whose right hand sides are quadratic in $(\hat\psi,\hat\varphi)$. The
	characteristic equation of the linear part is exactly $\Lambda(\lambda)=0$, so
	the equilibrium $(\hat\psi,\hat\psi',\hat\varphi,\hat\varphi')=(0,0,0,0)$ of
	the associated four-dimensional first order system is hyperbolic, with a
	two-dimensional unstable manifold (corresponding to $\lambda_1,\lambda_2>0$)
	and a two-dimensional stable manifold. Since $(\hat\psi,\hat\varphi)\to0$ as
	$x\to-\infty$, the solution lies on the unstable manifold, and by the theory
	of asymptotic integration \cite[Chapter 3, Theorem 8.1]{MR69338} (see also
	\cite{MR658490} for the degenerate cases) there are constants
	$\alpha_1,\alpha_2$, not both $0$, and eigenvectors $\xi^{(j)}$ of the linear
	part associated with $\lambda_j$, such that
	\[
	(\hat\psi,\hat\varphi)(x)=\alpha_1\xi^{(1)}e^{\lambda_1x}(1+o(1))
	\quad\text{or}\quad
	(\hat\psi,\hat\varphi)(x)=\alpha_2\xi^{(2)}e^{\lambda_2x}(1+o(1))
	\quad\text{ as }x\to-\infty,
	\]
	according to whether $\alpha_1\ne0$ or $\alpha_1=0$. From the first line of
	the linear system, an eigenvector associated with $\lambda$ satisfies
	$p(\lambda)\xi_1+c_1u^*\xi_2=0$, i.e. $\xi_2=-p(\lambda)\xi_1/(c_1u^*)$.
	Since $\lambda_1<\alpha_+$ we have $p(\lambda_1)<0$, so $\xi^{(1)}$ has two
	components of the same sign; since $\lambda_2>\alpha_+$ we have
	$p(\lambda_2)>0$, so $\xi^{(2)}$ has components of opposite signs. As
	$\hat\psi>0$ and $\hat\varphi>0$, the second alternative is impossible.
	Therefore $\alpha_1\ne0$ and \eqref{asym} holds with $k_1,k_2>0$.
\end{proof}

\begin{lemma}\label{lm:mono}
	$\psi_L'<0$ on $(-\infty,L]$ and $\varphi_L'>0$ on $\mathbb{R}$.
\end{lemma}

\begin{proof}
	For $s>0$ set $\psi^s(x):=\psi_L(x+s)$ and $\varphi^s(x):=\varphi_L(x+s)$.
	Since $A$ is nonincreasing, $A(x+s)\le A(x)$, so on $(-\infty,L-s)$
	\[
	-c(\psi^{s})' - d_1(\psi^{s})'' = \big(A(x+s)-b_1\psi^{s}-c_1\varphi^{s}\big)\psi^{s}
	\le \big(A(x)-b_1\psi^{s}-c_1\varphi^{s}\big)\psi^{s},
	\]
	while $-c(\varphi^{s})'-d_2(\varphi^{s})''=(a_2-b_2\psi^{s}-c_2\varphi^{s})\varphi^{s}$
	on $\mathbb{R}$. In the cooperative variables \eqref{cooper}, this says that
	$\Theta^{s}:=(\psi^{s},\frac{a_2}{c_2}-\varphi^{s})$ is a lower solution of
	\eqref{eqcoop} on $(-\infty,L-s)\times \mathbb{R}$, while $\Theta$ is a solution.

	We use the sliding method. Let
	\[
	S:=\Big\{s>0:\ \psi^{\sigma}\le\psi_L \text{ on }(-\infty,L-\sigma]\ \text{ and }\
	\varphi^{\sigma}\ge\varphi_L \text{ on }\mathbb{R},\ \ \forall\sigma\ge s\Big\}.
	\]

	\emph{Strict inequalities near $-\infty$.} By \cref{lm:asym}, for every
	$\varepsilon\in(0,1)$ there exists $X_\varepsilon<0$ such that
	\begin{equation}\label{eq:two_sided}
		\begin{aligned}
			(1-\varepsilon)k_1e^{\lambda_1y}&\le u^*-\psi_L(y)\le (1+\varepsilon)k_1e^{\lambda_1y},\\
			(1-\varepsilon)k_2e^{\lambda_1y}&\le \varphi_L(y)-v^*\le (1+\varepsilon)k_2e^{\lambda_1y},
		\end{aligned}
		\qquad (y\le X_\varepsilon).
	\end{equation}
	If $\sigma>0$ and $\varepsilon$ satisfy $(1-\varepsilon)e^{\lambda_1\sigma}>1+\varepsilon$,
	then for $x\le X_\varepsilon-\sigma$ both $x$ and $x+\sigma$ lie below
	$X_\varepsilon$, and \eqref{eq:two_sided} gives
	\[
	u^*-\psi^{\sigma}(x)\ \ge\ (1-\varepsilon)k_1e^{\lambda_1\sigma}e^{\lambda_1x}
	\ >\ (1+\varepsilon)k_1e^{\lambda_1x}\ \ge\ u^*-\psi_L(x)
	\]
	and, in the same way, $\varphi^{\sigma}(x)-v^*>\varphi_L(x)-v^*$. Hence
	\begin{equation}\label{eq:strict_minus}
		\psi^{\sigma}<\psi_L \quad\text{and}\quad \varphi^{\sigma}>\varphi_L
		\qquad\text{ on } (-\infty,X_\varepsilon-\sigma]
	\end{equation}
	whenever $(1-\varepsilon)e^{\lambda_1\sigma}>1+\varepsilon$. Note that, given
	any $\sigma_*>0$, this condition holds for all $\sigma\ge\sigma_*$ as soon as
	$\varepsilon$ is small enough.

	\emph{$S\ne\emptyset$.} Take $\varepsilon=\frac12$, write $X:=X_{1/2}$ and let
	$\sigma\ge\sigma_1:=\lambda_1^{-1}\ln5$, so that \eqref{eq:strict_minus} holds
	on $(-\infty,X-\sigma]$. We cover the remaining ranges using that
	$\psi_L<u^*$ on $(-\infty,L]$, that $\varphi_L>v^*$ on $\mathbb{R}$, and that, by
	\eqref{s0} and \cref{lm:phi_infty}, $\varphi_L\le\Phi_{c_0}(L)<\frac{a_2}{c_2}$
	on $(-\infty,L]$ while $\varphi_L$ is increasing on $(L,+\infty)$ with
	$\varphi_L(+\infty)=\frac{a_2}{c_2}$. Fix $Z>L$ with
	$\varphi_L>\Phi_{c_0}(L)$ on $[Z,+\infty)$ and set
	$M_0:=\max_{[X,L]}\psi_L<u^*$ and $m_0:=\min_{[X,Z]}(\varphi_L-v^*)>0$.
	\begin{itemize}
		\item For $x\in[X-\sigma,\,L-\sigma]$ we have $x+\sigma\in[X,L]$, hence
		$\psi^\sigma(x)\le M_0$; since $x\le L-\sigma$ and $\psi_L(-\infty)=u^*>M_0$,
		we get $\psi^\sigma(x)<\psi_L(x)$ for $\sigma$ large. As
		$\psi^\sigma\equiv0$ on $[L-\sigma,+\infty)$, the $\psi$-inequality holds on
		all of $(-\infty,L-\sigma]$.
		\item For $x\in[X-\sigma,\,Z-\sigma]$ we have $x+\sigma\in[X,Z]$, hence
		$\varphi^\sigma(x)-v^*\ge m_0$, while
		$\varphi_L(x)-v^*\le\frac32k_2e^{\lambda_1(Z-\sigma)}<m_0$ for $\sigma$ large.
		\item For $x\in[Z-\sigma,\,L]$ we have $x+\sigma\ge Z$, hence
		$\varphi^\sigma(x)>\Phi_{c_0}(L)\ge\varphi_L(x)$.
		\item For $x\ge L$ we have $\varphi^\sigma(x)=\varphi_L(x+\sigma)>\varphi_L(x)$,
		since $\varphi_L$ is increasing on $(L,+\infty)$.
	\end{itemize}
	Hence $S\supset[s_1,+\infty)$ for $s_1$ large.

	\emph{$\bar s:=\inf S=0$.} By continuity, $\bar s\in S$. Suppose
	$\bar s>0$, and fix $\varepsilon\in(0,1)$ so small that
	$(1-\varepsilon)e^{\lambda_1\bar s/2}>1+\varepsilon$; write
	$X_*:=X_\varepsilon-\bar s$, so that, by \eqref{eq:strict_minus},
	\begin{equation}\label{eq:unif_minus}
		\psi^{\sigma}<\psi_L \quad\text{and}\quad \varphi^{\sigma}>\varphi_L
		\qquad\text{ on } (-\infty,X_*] \text{ for every } \sigma\in[\bar s/2,\bar s].
	\end{equation}
	Set $P:=\psi_L-\psi^{\bar s}\ge0$ on $(-\infty,L-\bar s]$ and
	$Q:=\varphi^{\bar s}-\varphi_L\ge0$ on $\mathbb{R}$. By the computation above and the
	mean value theorem, $(P,Q)$ satisfies a linear cooperative system
	\[
	d_1P''+cP'+\gamma_{11}P+\gamma_{12}Q\le 0,\qquad
	d_2Q''+cQ'+\gamma_{21}P+\gamma_{22}Q\le 0
	\]
	with bounded coefficients and $\gamma_{12},\gamma_{21}\ge0$. Since
	$P(L-\bar s)=\psi_L(L-\bar s)>0$ and, by \eqref{eq:unif_minus},
	$P>0$ and $Q>0$ on $(-\infty,X_*]$, the strong maximum principle for
	cooperative systems yields
	\[
	P>0 \text{ on }(-\infty,L-\bar s],\qquad Q>0 \text{ on } \mathbb{R} .
	\]
	Choose $Z_*>L$ with $\varphi_L>\Phi_{c_0}(L)$ on $[Z_*,+\infty)$. On the
	compact sets $[X_*,L-\bar s]$ and $[X_*,Z_*]$ the continuous functions $P$
	and $Q$ are bounded below by a positive constant; since moreover
	$\psi_L>0$ on $(-\infty,L)$ and $\psi^\sigma(L-\sigma)=0$, a standard
	compactness argument provides $\delta\in(0,\bar s/2)$ such that, for every
	$\sigma\in[\bar s-\delta,\bar s]$,
	\[
	\psi^{\sigma}<\psi_L \ \text{ on }[X_*,L-\sigma]
	\qquad\text{and}\qquad
	\varphi^{\sigma}>\varphi_L \ \text{ on }[X_*,Z_*].
	\]
	Outside these sets the required inequalities hold by \eqref{eq:unif_minus}
	(for $x\le X_*$) and, for $x\ge Z_*$, because
	$\varphi^\sigma(x)=\varphi_L(x+\sigma)>\varphi_L(x)$ by the monotonicity of
	$\varphi_L$ on $(L,+\infty)$. Hence $\bar s-\delta\in S$, contradicting the
	definition of $\bar s$.

	Therefore $\bar s=0$, i.e. $\psi_L(x+s)\le\psi_L(x)$ and
	$\varphi_L(x+s)\ge\varphi_L(x)$ for all $s>0$, so $\psi_L$ is nonincreasing
	and $\varphi_L$ is nondecreasing. Applying the strong maximum principle to
	$\psi_L'$ and $\varphi_L'$, which satisfy linear equations obtained by
	differentiating \eqref{eqlm1} (recall $A$ is Lipschitz and nonincreasing, so
	$A'\le0$ a.e.), we conclude that $\psi_L'<0$ on $(-\infty,L]$ and
	$\varphi_L'>0$ on $\mathbb{R}$.
\end{proof}

\begin{lemma}\label{lm:uniq}
	The solution of \eqref{eqlm1} is unique.
\end{lemma}

\begin{proof}
	Let $(\psi_1,\varphi_1)$ and $(\psi_2,\varphi_2)$ be two solutions of
	\eqref{eqlm1}. By \cref{lm:asym}, there are $\lambda_1>0$ and positive
	constants $k_1^{(i)},k_2^{(i)}$ with
	\[
	u^*-\psi_i(x)=\big(k_1^{(i)}+o(1)\big)e^{\lambda_1x},
	\qquad
	\varphi_i(x)-v^*=\big(k_2^{(i)}+o(1)\big)e^{\lambda_1x}\qquad(x\to-\infty),\ i=1,2 .
	\]
	For $\xi\ge0$ set $\psi_2^\xi(x):=\psi_2(x+\xi)$ and
	$\varphi_2^\xi(x):=\varphi_2(x+\xi)$. Exactly as in the proof of
	\cref{lm:mono}, $(\psi_2^\xi,\varphi_2^\xi)$ is a lower solution of
	\eqref{eqlm1} on $(-\infty,L-\xi]$, and the two-sided bounds
	\eqref{eq:two_sided}, now written with the constants $k_j^{(i)}$, give: if
	$\varepsilon\in(0,1)$ and $\xi>0$ satisfy
	$(1-\varepsilon)k_j^{(2)}e^{\lambda_1\xi}>(1+\varepsilon)k_j^{(1)}$ for
	$j=1,2$, then $\psi_2^\xi<\psi_1$ and $\varphi_2^\xi>\varphi_1$ on a half-line
	$(-\infty,X]$; and for any $\xi_*>0$ this holds for all $\xi\ge\xi_*$ provided
	$\varepsilon$ is small enough.

	The only point that requires an argument different from \cref{lm:mono} is the
	behaviour at $+\infty$, where $\varphi_1$ and $\varphi_2^\xi$ are two
	\emph{different} functions with the same limit $\frac{a_2}{c_2}$. For
	$x\ge L$ we have $\psi_1(x)=\psi_2^\xi(x)=0$, so both $\varphi_1$ and
	$\varphi_2^\xi$ solve the scalar equation
	$-c\theta'-d_2\theta''=(a_2-c_2\theta)\theta$ there, and
	$W:=\varphi_2^\xi-\varphi_1$ satisfies
	\[
	d_2W''+cW'+qW=0 \quad\text{ on }(L,+\infty),
	\qquad q:=a_2-c_2\big(\varphi_2^\xi+\varphi_1\big).
	\]
	Fix $\epsilon\in\big(0,\frac{a_2}{2c_2}\big)$ and choose $R>L$ so large that
	$\varphi_1>\frac{a_2}{c_2}-\epsilon$ and $\varphi_2>\frac{a_2}{c_2}-\epsilon$
	on $[R,+\infty)$; since $\xi\ge0$ and $\varphi_2$ is increasing on
	$(L,+\infty)$, the same bound holds for $\varphi_2^\xi$ on $[R,+\infty)$,
	whence $q<-a_2+2c_2\epsilon<0$ there, uniformly in $\xi\ge0$. As
	$W(+\infty)=0$, a negative interior minimum of $W$ on $[R,+\infty)$ is
	impossible (at such a point $W''\ge0$, $W'=0$ and $qW>0$), so
	\begin{equation}\label{eq:tail}
		\inf_{[R,+\infty)}\big(\varphi_2^\xi-\varphi_1\big)\ \ge\
		\min\big\{0,\ \varphi_2^\xi(R)-\varphi_1(R)\big\} .
	\end{equation}
	Consequently it suffices to control $\varphi_2^\xi-\varphi_1$ on
	$(-\infty,R]$, and the sliding argument of \cref{lm:mono} applies verbatim
	with $Z$ (resp. $Z_*$) replaced by $R$: the set
	\[
	\mathcal S=\big\{\xi\ge0:\ \psi_2^\xi\le\psi_1 \text{ on }(-\infty,L-\xi],\
	\varphi_2^\xi\ge\varphi_1 \text{ on }\mathbb{R}\big\}
	\]
	is nonempty and $\bar\xi:=\inf\mathcal S=0$. Hence $\psi_2\le\psi_1$ and
	$\varphi_2\ge\varphi_1$. Exchanging the roles of the two solutions gives the
	reverse inequalities, so $(\psi_1,\varphi_1)=(\psi_2,\varphi_2)$.
\end{proof}


\subsection{Proof of \texorpdfstring{\cref{th:semi-wave}}{the semi-wave theorem}(ii) and (iii)}

\begin{proof}[Proof of \cref{th:semi-wave}(ii)]
	Let $0\le L_1<L_2$ and set $\psi_2(x):=\psi_{L_2}(x+L_2-L_1)$,
	$\varphi_2(x):=\varphi_{L_2}(x+L_2-L_1)$. Then $(\psi_2,\varphi_2)$ satisfies
	\[
	\begin{cases}
		-c \psi_2' - d_1 \psi_2'' = \big(A(x + L_2 - L_1) - b_1\psi_2 - c_1\varphi_2\big)\psi_2, & -\infty < x < L_1,\\
		-c \varphi_2' - d_2 \varphi_2'' = (a_2 - b_2\psi_2 - c_2\varphi_2)\varphi_2, & -\infty < x < +\infty,\\
		\psi_2(-\infty) = u^*, ~ \psi_2(x)=0 \text{ for } x\ge L_1,\\
		\varphi_2(-\infty) = v^*, ~ \varphi_2(+\infty)=\frac{a_2}{c_2} .
	\end{cases}
	\]
	Since $A(x+L_2-L_1)\le A(x)$, the pair $(\psi_2,\varphi_2)$ is a lower
	solution of \eqref{eqlm1} with $L=L_1$. Both pairs have the same limits at
	$-\infty$, both $\psi$-components vanish at $L_1$, and by \cref{lm:asym}
	(applied to $\psi_{L_2},\varphi_{L_2}$ and then translated) the pair
	$(\psi_2,\varphi_2)$ obeys the asymptotics \eqref{asym} with the same exponent
	$\lambda_1$. The sliding argument of \cref{lm:uniq} therefore applies verbatim
	and gives
	\[
	\psi_{L_1}\ge\psi_2 \text{ on }(-\infty,L_1],
	\qquad \varphi_{L_1}\le\varphi_2 \text{ on }\mathbb{R} .
	\]
	Both $\psi$-components vanish at $x=L_1$; moreover
	$\psi_{L_1}\not\equiv\psi_2$, since $A$ is strictly decreasing on $[0,l_0]$
	and $L_2>L_1$, so that $A(x+L_2-L_1)<A(x)$ on a nonempty open subset of
	$(-\infty,L_1)$. The strong maximum principle for cooperative systems then
	yields $\psi_{L_1}>\psi_2$ on $(-\infty,L_1)$, and the Hopf boundary lemma
	gives
	\[
	\psi_{L_1}'(L_1) < \psi_2'(L_1) = \psi_{L_2}'(L_2),
	\]
	that is, $L\mapsto\psi_L'(L)$ is strictly increasing.

	We next show $\psi_L'(L)\to0$ as $L\to+\infty$. Put $w_L(x):=\psi_L(x+L)$ for
	$x\le0$, so that $w_L(0)=0$, $0<w_L\le u^*$ and, by \eqref{s0},
	\[
	-cw_L'-d_1w_L''=\big(A(x+L)-b_1w_L-c_1\varphi_L(x+L)\big)w_L
	\le\big(A(x+L)-c_1v^*\big)w_L .
	\]
	For $x\in I_L:=(l_0-L,0)$ we have $x+L>l_0$, hence $A(x+L)=a_0$ and,
	with $\kappa:=c_1v^*-a_0>0$,
	\[
	d_1w_L''+cw_L'-\kappa w_L\ \ge\ 0 \qquad\text{ in } I_L .
	\]
	Let $\beta:=\frac{-c-\sqrt{c^2+4d_1\kappa}}{2d_1}<0$, so that
	$d_1\beta^2+c\beta-\kappa=0$, and set
	$W(x):=u^* e^{\beta(x-l_0+L)}$. Then $d_1W''+cW'-\kappa W=0$,
	$W(l_0-L)=u^*\ge w_L(l_0-L)$ and $W(0)>0=w_L(0)$. The maximum principle
	applied to $w_L-W$ on $I_L$ gives $w_L\le W$ on $I_L$, hence
	\[
	\|w_L\|_{L^\infty([-2,0])}\ \le\ u^* e^{\beta(L-l_0-2)}\ \longrightarrow\ 0
	\qquad (L\to+\infty).
	\]
	Since the coefficients of the equation for $w_L$ are bounded uniformly in
	$L$, elliptic estimates on $[-2,0]$ give
	$\|w_L\|_{C^1([-1,0])}\to0$, and in particular
	$\psi_L'(L)=w_L'(0)\to0$.
\end{proof}

\begin{proof}[Proof of \cref{th:semi-wave}(iii)]
	Define $f(L):=-\mu(A(L))\psi_L'(L)$ for $L\ge0$. By part (ii), $L\mapsto
	-\psi_L'(L)$ is positive, strictly decreasing and tends to $0$; since $A$ is
	nonincreasing and $\mu$ is nondecreasing with $\mu(a_0)>0$, the function
	$L\mapsto\mu(A(L))$ is positive and nonincreasing. Hence $f$ is positive,
	strictly decreasing and $f(+\infty)=0$. Continuity of $f$ follows from the
	uniqueness in part (i) together with elliptic estimates and a compactness
	argument.

	For $L=0$ we have $A\equiv a_1$ on $(-\infty,0]$, so \eqref{eqlm1} coincides
	with \eqref{semiwave0}; by the uniqueness in \cref{th:semi-wave0},
	$(\psi_0,\varphi_0)\equiv(\Psi_c,\Phi_c)$ and therefore
	\[
	f(0)=-\mu(a_1)\Psi_c'(0).
	\]
	Consider $G(c):=-\mu(a_1)\Psi_c'(0)-c$ on $[0,c_0]$. It is continuous by
	\cref{th:semi-wave0}, $G(0)=-\mu(a_1)\Psi_0'(0)>0$, $G(c_0)=0$ by \eqref{c0},
	and $c_0$ is the \emph{only} zero of $G$ in $[0,c^*)$. Hence $G>0$ on
	$[0,c_0)$, i.e.
	\[
	f(0)\ \ge\ c \quad\text{ for } 0<c\le c_0,\qquad\text{ with equality iff } c=c_0 .
	\]
	Since $f$ is continuous, strictly decreasing, $f(0)\ge c$ and $f(+\infty)=0$,
	the intermediate value theorem provides a unique $L_0\ge0$ with $f(L_0)=c$,
	and $L_0=0$ if and only if $f(0)=c$, i.e. if and only if $c=c_0$. In that case
	$(\psi_0,\varphi_0)\equiv(\Psi_{c_0},\Phi_{c_0})$.
\end{proof}

\begin{remark}\label{rm:perturbed}
	In \cref{sec_6} we shall apply \cref{th:semi-wave} to perturbed problems in
	which $(a_1,a_2)$ is replaced by $(a_1\pm2\delta,a_2\mp\delta)$ and $A$ by
	$A\pm2\delta$. For $\delta>0$ small all the structural assumptions are
	preserved: \eqref{weak_competition} remains valid by continuity, the
	perturbed $A$ still satisfies \eqref{A} with $a_0\pm2\delta<0<a_1\pm2\delta$,
	and $\mu$ is defined on the perturbed range thanks to the extension of $\mu$
	introduced after \eqref{mu}. Denoting by $c_{0,\delta}$ the corresponding
	critical speed, the continuous dependence of $c_0$ on the parameters
	(\cref{th:semi-wave0}) gives $c_{0,\delta}\to c_0$ as $\delta\to0$.
\end{remark}

\section{The spreading speed: proof of \texorpdfstring{\cref{th:spreading_speed}}{the spreading speed theorem}}\label{sec_6}

Throughout this section we assume \eqref{weak_competition}, \eqref{initial} and
that spreading occurs, i.e. $h_\infty=+\infty$; by \cref{th1:lm2},
\begin{equation}\label{eq:spread_conv}
	\lim_{t\to+\infty}\big(u(t,x),v(t,x)\big)=(u^*,v^*)
	\qquad\text{ locally uniformly in }x\in[0,\infty).
\end{equation}
We distinguish the three cases $c<c_0$, $c=c_0$ and $c>c_0$. Several comparison
arguments below start at a positive time $T$; this is legitimate because
\cref{th7} holds with $0$ replaced by any initial time $T\ge0$, the data
$(u_0,v_0,h_0)$ being replaced by $(u(T,\cdot),v(T,\cdot),h(T))$. We always keep
the original time variable, so that the environment felt by the comparison
functions is exactly $A(\cdot-ct)$ and no relabelling of $A$ is needed.

\subsection{The case \texorpdfstring{$c<c_0$}{c<c0}}

\begin{lemma}\label{lem:limsup}
	Suppose $c < c_0$. Then $\displaystyle\limsup_{t \to +\infty} \frac{h(t)}{t} \le c$.
\end{lemma}

\begin{proof}
	For $\delta>0$ set
	\[
	u^*_\delta:=\frac{(a_1+2\delta)c_2-(a_2-\delta)c_1}{b_1c_2-b_2c_1},
	\qquad
	v^*_\delta:=\frac{b_1(a_2-\delta)-(a_1+2\delta)b_2}{b_1c_2-b_2c_1},
	\]
	which are the coexistence states of the system obtained from \eqref{main} by
	replacing $A$ with $A+2\delta$ and $a_2$ with $a_2-\delta$. Clearly
	$\delta\mapsto u^*_\delta$ is increasing and $\delta\mapsto v^*_\delta$ is
	decreasing, so that
	\begin{equation}\label{eq:order_delta}
		u^* < u^*_{\delta/2} < u^*_\delta,
		\qquad
		v^* > v^*_{\delta/2} > v^*_\delta .
	\end{equation}
	By \cref{rm:perturbed}, for $\delta>0$ small all the structural assumptions
	are preserved and the corresponding critical speed $c_{0,\delta}$ satisfies
	$c_{0,\delta}\to c_0$ as $\delta\to0$. Since $c<c_0$, we may and do fix
	$\delta>0$ so small that $c<c_{0,\delta}$. By \cref{th:semi-wave} applied to
	the perturbed problem, there exist $L_\delta\ge0$ and a unique solution
	$(\psi_\delta,\varphi_\delta)$ of
	\begin{equation}\label{aux}
		\begin{cases}
			-c \psi' - d_1 \psi'' = (A(x) +2 \delta - b_1\psi - c_1\varphi)\psi, & -\infty < x < L_\delta,\\
			-c \varphi' - d_2 \varphi'' = (a_2 - \delta - b_2\psi - c_2\varphi)\varphi, & -\infty < x < +\infty,\\
			\psi(-\infty) = u^*_\delta, ~ \psi(x)=0 \text{ for } x\ge L_\delta, ~ \psi'(x)<0 \text{ for } x\le L_\delta,\\
			\varphi(-\infty) = v^*_\delta, ~ \varphi(+\infty)=\frac{a_2 - \delta}{c_2}, ~ \varphi'(x)>0 \text{ for } x\in \mathbb{R},
		\end{cases}
	\end{equation}
	such that
	\begin{equation}\label{semi-wave-c}
		c = -\mu\big(A(L_\delta) +2 \delta\big) \psi'_{\delta}(L_\delta).
	\end{equation}
	Note that $0<\psi_\delta<u^*_\delta$ on $(-\infty,L_\delta)$ and
	$v^*_\delta<\varphi_\delta<\frac{a_2-\delta}{c_2}$ on $\mathbb{R}$.

	\textit{Step 1.} There exists $T_1>0$ such that
	\begin{equation}\label{eq:step1}
		u(t,x) \le u^*_{\delta/2},\qquad v(t,x) \ge v^*_{\delta/2}
		\qquad (x\ge0,\ t\ge T_1).
	\end{equation}
	Indeed, this is \eqref{eq:uv_global} of \cref{lem:vlow} combined with
	\eqref{eq:order_delta}.

	\textit{Step 2.} For every $T_0>0$ there exist $T>T_0$ and $S_0>0$ such that
	\begin{equation}\label{eq:step2}
		v(T, x) \ge \frac{a_2-\delta}{c_2} \qquad\text{ for all } x\ge S_0 .
	\end{equation}
	To see this, recall $\sigma_0^{\max}=\inf_{x\ge0}v_0(x)>0$ from \eqref{v0inf}
	and consider
	\begin{equation}\label{aux1}
		\begin{cases}
			w_t - d_2w_{xx} = (a_2 - c_2w)w, & t > 0, ~ x > 0,\\
			w(t, 0) = 0, & t > 0,\\
			w(0,x) = \min\{\sigma_0^{\max},a_2/c_2\}, & x > 0.
		\end{cases}
	\end{equation}
	It is classical that $w(t,\cdot)\to w_*$ locally uniformly on $[0,\infty)$,
	where $w_*$ is the unique positive solution of $-d_2w_*''=(a_2-c_2w_*)w_*$ on
	$(0,\infty)$ with $w_*(0)=0$, and that $w_*'>0$, $w_*(+\infty)=a_2/c_2$.
	Choose $S_1>0$ and $T>T_0$ with
	$w(t,S_1)\ge w_*(S_1)-\frac{\delta}{2c_2}\ge\frac{a_2-\delta}{c_2}$ for
	$t\ge T$. Applying the maximum principle to the equation satisfied by $w_x$
	gives $w_x\ge0$, hence
	\begin{equation}\label{eq:w_lower}
		w(t,x)\ge\frac{a_2-\delta}{c_2}\qquad (t\ge T,\ x\ge S_1).
	\end{equation}
	Set $S_2:=h(T)$ and $\tilde w(t,x):=w(t,x-S_2)$. Since $h(t)\le S_2$ for
	$t\in[0,T]$, the function $v$ satisfies
	$v_t-d_2v_{xx}=(a_2-c_2v)v$ for $0<t\le T$, $x>S_2$, together with
	$\tilde w(t,S_2)=0<v(t,S_2)$ and
	$\tilde w(0,x)\le\sigma_0^{\max}\le v(0,x)$ for $x>S_2$; the comparison
	principle yields $v(t,x)\ge w(t,x-S_2)$ for $0<t\le T$ and $x>S_2$. With
	$S_0:=S_1+S_2$ and \eqref{eq:w_lower} we obtain \eqref{eq:step2}.

	\textit{Step 3: construction of an upper solution.} Apply Step 2 with
	$T_0:=T_1$ to obtain $T>T_1$ and $S_0>0$ as in \eqref{eq:step2}, and set
	$S:=S_0+h(T)$. Since $\psi_\delta(-\infty)=u^*_\delta>u^*_{\delta/2}$ and
	$\varphi_\delta(-\infty)=v^*_\delta<v^*_{\delta/2}$, we may choose $R>0$ so
	large that
	\begin{equation}\label{eq:choiceR}
		cT+L_\delta+R>h(T),
		\qquad
		\psi_{\delta}(x -cT-R) > u^*_{\delta/2}
		\quad\text{and}\quad
		\varphi_{\delta}(x -cT- R) < v^*_{\delta/2}
		\qquad\text{ for } x\in[0,S].
	\end{equation}
	Define, for $t\ge T$,
	\[
	\overline{h}(t) := c t + L_\delta + R,\qquad
	\overline{u}(t,x) := \psi_{\delta}(x-ct -R),\qquad
	\underline{v}(t,x) := \varphi_{\delta}(x-ct-R).
	\]
	We claim that $(\overline u,\underline v,\overline h)$ is an upper solution of
	\eqref{main} on $[T,+\infty)$. Note that we keep the original time variable
	here, so that the environment felt by the comparison functions is exactly
	$A(\cdot-ct)$ and no relabelling of $A$ is needed. Indeed, since $A$ is nonincreasing and
	$R>0$, we have $A(x-ct-R)+2\delta\ge A(x-ct)$, so
	\begin{align*}
		\overline{u}_t - d_1 \overline{u}_{xx}
		&= -c\psi'_\delta-d_1\psi''_\delta
		= \big(A(x - c t - R) + 2\delta - b_1 \overline u - c_1 \underline v\big) \overline u
		\ \ge\ \big(A(x - ct) - b_1 \overline{u} - c_1 \underline{v}\big) \overline{u},\\
		\underline{v}_t - d_2 \underline{v}_{xx}
		&= -c\varphi'_\delta-d_2\varphi''_\delta
		= \big(a_2 - \delta - b_2\overline u - c_2\underline v\big) \underline v
		\ \le\ \big(a_2 - b_2\overline{u} - c_2\underline{v}\big)\underline{v},
	\end{align*}
	and
	\[
	\overline{u}_x(t, 0)=\psi'_\delta(-ct-R) \le 0 \le \varphi'_\delta(-ct-R)=\underline{v}_x(t, 0),
	\qquad
	\overline{u}(t, x) = 0 \ \text{ for } x\ge \overline{h}(t) ,
	\]
	so the first alternative at $x=0$ in \cref{th7}(a) holds. For the free boundary
	condition, note $\overline h(t)-ct-R=L_\delta$, hence
	$\overline u_x(t,\overline h(t))=\psi'_\delta(L_\delta)$, and since
	$A(L_\delta)+2\delta> A(L_\delta)\ge A(L_\delta+R)=A(\overline h(t)-ct)$ and
	$\mu$ is nondecreasing, \eqref{semi-wave-c} gives
	\[
	\overline{h}'(t) = c = -\mu\big(A(L_\delta)+2 \delta\big) \psi_{\delta}'(L_\delta)
	\ \ge\ -\mu\big(A(\overline h(t)-ct)\big) \overline u_x(t,\overline h(t)) .
	\]
	Finally, by \eqref{eq:choiceR}, \eqref{eq:step1} and \eqref{eq:step2},
	\[
	\overline h(T)=cT+L_\delta+R>h(T),\qquad
	\overline{u}(T, x) = \psi_\delta(x-cT-R)>u^*_{\delta/2}\ge u(T,x)\quad (0\le x\le h(T)),
	\]
	\[
	\underline{v}(T, x) = \varphi_\delta(x-cT-R)
	\begin{cases}
		<v^*_{\delta/2}\le v(T,x), & 0\le x\le S,\\[2pt]
		<\frac{a_2-\delta}{c_2}\le v(T,x), & x> S\ (\ge S_0).
	\end{cases}
	\]
	By \cref{th7}(a), $h(t)\le\overline{h}(t)=ct+L_\delta+R$ for all $t\ge T$,
	whence $\limsup_{t\to+\infty}h(t)/t\le c$.
\end{proof}

\begin{lemma}\label{lem:liminf}
	Suppose $c < c_0$. Then $\displaystyle\liminf_{t \to +\infty} \frac{h(t)}{t} \ge c$.
\end{lemma}

\begin{proof}
	For $\delta>0$ set
	\[
	u^*_\delta:=\frac{(a_1-2\delta)c_2-(a_2+2\delta)c_1}{b_1c_2-b_2c_1},
	\qquad
	v^*_\delta:=\frac{b_1(a_2+2\delta)-(a_1-2\delta)b_2}{b_1c_2-b_2c_1},
	\]
	so that $u^*_\delta<u^*$ and $v^*_\delta>v^*$. As in \cref{lem:limsup},
	\cref{rm:perturbed} allows us to fix $\delta>0$ so small that
	\eqref{weak_competition} still holds for the perturbed parameters and $c$ is
	below the corresponding critical speed. By \cref{th:semi-wave} there exist
	$L_\delta\ge0$ and a unique solution $(\psi_\delta,\varphi_\delta)$ of
	\begin{equation}\label{aux2}
		\begin{cases}
			-c \psi' - d_1 \psi'' = (A(x) -2\delta - b_1\psi - c_1\varphi)\psi, & -\infty < x < L_\delta,\\
			-c \varphi' - d_2 \varphi'' = (a_2 + 2\delta - b_2\psi - c_2\varphi)\varphi, & -\infty < x < +\infty,\\
			\psi(-\infty) = u^*_\delta, ~ \psi(x)=0 \text{ for } x\ge L_\delta, ~ \psi'(x)<0 \text{ for } x\le L_\delta,\\
			\varphi(-\infty) = v^*_\delta, ~ \varphi(+\infty)=\frac{a_2 + 2\delta}{c_2}, ~ \varphi'(x)>0 \text{ for } x\in \mathbb{R},
		\end{cases}
	\end{equation}
	with
	\begin{equation}\label{eq:sw2}
		c = -\mu\big(A(L_\delta) - 2\delta\big)\psi_{\delta}'(L_\delta).
	\end{equation}
	In particular $0<\psi_\delta<u^*_\delta$ on $(-\infty,L_\delta)$ and
	$v^*_\delta<\varphi_\delta<\frac{a_2+2\delta}{c_2}$ on $\mathbb{R}$.

	Let $V$ be the solution of $V'=(a_2-c_2V)V$, $V(0)=\|v_0\|_\infty$; comparison
	gives $v(t,x)\le V(t)$ for all $t>0$, $x\ge0$, and since
	$V(t)\to a_2/c_2$ there is $T_2>0$ with
	\begin{equation}\label{eq:v}
		v(t,x) < \frac{a_2 + \delta}{c_2} \qquad (t \ge T_2,\ x \ge 0).
	\end{equation}
	Since $\varphi_\delta$ is increasing with
	$\varphi_\delta(+\infty)=\frac{a_2+2\delta}{c_2}$, we may pick
	$R_1\ge L_\delta$ with
	\begin{equation}\label{eq:vp}
		\varphi_{\delta}(x) > \frac{a_2+\delta}{c_2} \quad\text{ for } x \ge R_1 .
	\end{equation}
	By \eqref{eq:spread_conv} and $h_\infty=+\infty$, and since $u^*_\delta<u^*$
	and $v^*_\delta>v^*$, there exists $T_0>T_2$ such that
	\begin{equation}\label{eq:uv}
		h(T_0) > L_\delta,\qquad
		u(t,x) > u^*_\delta,\qquad v(t,x) < v^*_\delta
		\qquad (t \ge T_0,\ x \in [0, R_1]).
	\end{equation}

	Define, for $t\ge T_0$,
	\[
	\underline{h}(t) := c(t-T_0) + L_\delta,\qquad
	\underline{u}(t,x) := \psi_{\delta}\big(x - c(t-T_0)\big),\qquad
	\overline{v}(t,x) := \varphi_{\delta}\big(x - c(t-T_0)\big).
	\]
	We claim that $(\underline u,\overline v,\underline h)$ is a lower solution of
	\eqref{main} with initial time $T_0$. First, since $A$ is nonincreasing and
	$cT_0>0$, we have $A\big(x-c(t-T_0)\big)\le A(x-ct)$, hence for
	$x\in[0,\underline h(t)]$,
	\begin{align*}
		\underline{u}_t - d_1 \underline{u}_{xx}
		&= -c \psi_{\delta}' - d_1 \psi_{\delta}''
		= \big(A(x -c(t-T_0)) -2 \delta - b_1 \underline{u} - c_1 \overline{v}\big) \underline{u}
		\ \le\ \big(A(x - ct) - b_1 \underline{u} - c_1 \overline{v}\big) \underline{u},\\
		\overline{v}_t - d_2 \overline{v}_{xx}
		&= -c\varphi_{\delta}' -d_2 \varphi_{\delta}''
		= \big(a_2 + 2\delta - b_2 \underline u - c_2 \overline v\big) \overline v
		\ \ge\ \big(a_2 - b_2\underline{u} - c_2 \overline{v}\big) \overline{v}.
	\end{align*}
	Next, $\underline u_x(t,\underline h(t))=\psi_\delta'(L_\delta)$ and
	$\underline h(t)-ct=L_\delta-cT_0<L_\delta$, so
	$A(\underline h(t)-ct)\ge A(L_\delta)>A(L_\delta)-2\delta$ and, $\mu$ being
	nondecreasing and $-\psi'_\delta(L_\delta)>0$, \eqref{eq:sw2} gives
	\[
	\underline{h}'(t) = c = -\mu\big(A(L_\delta) - 2 \delta \big)\psi_{\delta}'(L_\delta)
	\ \le\ -\mu\big(A(\underline{h}(t) - ct)\big)\underline{u}_x(t,\underline{h}(t)).
	\]
	At $x=0$ we use the \emph{second} alternative of \cref{th7}(b): by
	\eqref{eq:uv}, for $t\ge T_0$,
	\[
	\underline{u}(t,0)=\psi_\delta\big(-c(t-T_0)\big)<u^*_\delta<u(t,0),
	\qquad
	\overline{v}(t,0)=\varphi_\delta\big(-c(t-T_0)\big)>v^*_\delta>v(t,0).
	\]
	Finally, at $t=T_0$ we have $\underline h(T_0)=L_\delta<h(T_0)$,
	\[
	\underline{u}(T_0,x) = \psi_{\delta}(x) < u^*_\delta < u(T_0,x)
	\qquad (0\le x \le L_\delta),
	\]
	and, using \eqref{eq:uv} for $x\in[0,R_1]$ and \eqref{eq:v}--\eqref{eq:vp} for
	$x>R_1$,
	\[
	\overline{v}(T_0,x) = \varphi_{\delta}(x)
	\begin{cases}
		>v^*_\delta> v(T_0,x), & 0\le x\le R_1,\\[2pt]
		>\frac{a_2+\delta}{c_2}>v(T_0,x), & x> R_1 .
	\end{cases}
	\]
	By \cref{th7}(b), $h(t)\ge\underline h(t)=c(t-T_0)+L_\delta$ for $t\ge T_0$,
	and therefore $\liminf_{t\to+\infty}h(t)/t\ge c$.
\end{proof}

\subsection{The case \texorpdfstring{$c\ge c_0$}{c>=c0}}

\begin{lemma}\label{lm:limsup_c0}
	Suppose $c\ge c_0$. Then $\displaystyle\limsup_{t\to+\infty}\frac{h(t)}{t}\le c_0$.
\end{lemma}

\begin{proof}
	As in Step~0 of the proof of \cref{th1:lm2}, $(u,v,h)$ is a lower solution of
	\eqref{main2}; let $(\overline u,\underline v,\overline h)$ be the solution of
	\eqref{main2} with the same initial data, so that $h\le\overline h$. Since
	$h_\infty=+\infty$ we get $\overline h(t)\to+\infty$, i.e. spreading occurs
	for \eqref{main2}, and \cite[Theorem 1.3]{MR3986328} gives
	$\overline h(t)/t\to c_0$. Hence $\limsup_{t\to+\infty}h(t)/t\le c_0$.
\end{proof}

\begin{lemma}\label{lm:eq}
	Suppose $c = c_0$. Then $\displaystyle\lim_{t\to+\infty} \frac{h(t)}{t} = c_0$.
\end{lemma}

\begin{proof}
	By \cref{lm:limsup_c0} it suffices to prove
	$\liminf_{t\to+\infty}h(t)/t\ge c_0$. For small $\delta>0$ put
	$c_\delta:=c_0-\delta>0$ and let $(u_\delta,v_\delta,h_\delta)$ be the
	solution of \eqref{main} with $c$ replaced by $c_\delta$ and with the same
	initial data. Since $c_\delta<c$ and $A$ is nonincreasing, we have
	$A(x-ct)\ge A(x-c_\delta t)$ and $\mu(A(h(t)-ct))\ge\mu(A(h(t)-c_\delta t))$
	for $t>0$; as $-u_x(t,h(t))>0$, this shows that $(u,v,h)$ is an upper solution
	of the $c_\delta$-problem in the sense of \cref{th7}(a) (first alternative at
	$x=0$). Hence
	\begin{equation}\label{eq:cdelta}
		h_\delta(t)\le h(t)\qquad (t\ge0).
	\end{equation}

	We next check that spreading occurs for the $c_\delta$-problem when $\delta$
	is small. Since $h_\infty=+\infty$, there is $T>0$ with $h(T)=R^*+1$. The
	solution of \eqref{main} depends continuously on the parameter $c$, uniformly
	on the fixed interval $[0,T]$ (this follows from the contraction mapping
	construction of \cref{th4}, whose constants depend continuously on $c$,
	together with uniqueness and a continuation argument on $[0,T]$, the solution
	being global with bounds uniform in $c$ by \cref{th6}), so we may fix
	$\delta>0$ small enough that $h_\delta(T)>R^*$. If vanishing occurred for the
	$c_\delta$-problem, \cref{th1:lm1} would give $h_{\delta,\infty}\le R^*$,
	contradicting $h_\delta(T)>R^*$. Hence spreading occurs for the
	$c_\delta$-problem.

	Since $c_\delta<c_0$, \cref{lem:liminf} applied to the $c_\delta$-problem
	yields $\liminf_{t\to+\infty}h_\delta(t)/t\ge c_\delta$, and by
	\eqref{eq:cdelta},
	\[
	\liminf_{t\to+\infty}\frac{h(t)}{t}\ \ge\ \liminf_{t\to+\infty}\frac{h_\delta(t)}{t}\ \ge\ c_0-\delta .
	\]
	Letting $\delta\to0$ completes the proof.
\end{proof}

\begin{lemma}\label{lm:g}
	Suppose $c > c_0$. Then $\displaystyle\lim_{t\to+\infty} \frac{h(t)}{t} = c_0$.
\end{lemma}

\begin{proof}
	By \cref{lm:limsup_c0}, $\limsup_{t\to+\infty}h(t)/t\le c_0<c$, so there
	exists $T_0>0$ with $h(t)<ct$ for all $t>T_0$. By \cref{lem:reduction},
	$(t,x)\mapsto(u(t+T_0,x),v(t+T_0,x),h(t+T_0))$ solves \eqref{main2} with data
	$(u(T_0,\cdot),v(T_0,\cdot),h(T_0))$, which satisfy \eqref{initial} (the
	condition $\liminf_{x\to\infty}v(T_0,x)>0$ following from \cref{lem:vlow}).
	Since $h_\infty=+\infty$, spreading occurs for that problem, and
	\cite[Theorem 1.3]{MR3986328} gives $h(t)/t\to c_0$.
\end{proof}

\begin{proof}[Proof of \cref{th:spreading_speed}]
	Combine \cref{lem:limsup}, \cref{lem:liminf}, \cref{lm:eq} and \cref{lm:g}.
\end{proof}

\section{Numerical simulations}\label{sec_7}

In this section we illustrate the theoretical results by numerical experiments.
Problem \eqref{main} is discretised by a finite difference scheme on the
straightened domain, combined with a Newton--Raphson iteration at each time step.

In all the figures below the parameters are fixed as
\[
d_1 = 0.02,\quad a_1 = 0.14,\quad b_1 = 0.4,\quad c_1 = 0.11,\quad
d_2 = 0.018,\quad a_2 = 0.18,\quad b_2 = 0.1,\quad c_2 = 0.2,
\]
\[
a_0 = -0.2,\quad l_0 = 1,\quad \mu(s) = 0.2 + 0.1s,\quad
A(\xi) = \begin{cases}
	a_1, & \xi \le 0, \\
	a_0 + \frac{1}{2} (a_1 - a_0)\big(\cos( \pi \xi / l_0 ) + 1\big), & 0 < \xi < l_0, \\
	a_0, & \xi \ge l_0,
\end{cases}
\]
and the initial data are
$u_0(x) = 0.2 \cos\big(\frac{\pi x}{2h_0}\big)$ and
$v_0(x) = 0.2 \cos\big(\frac{\pi x}{2}\big) + 0.3$, where $h_0$ is the initial
size of the invaded range. These values satisfy all our standing assumptions:
$\mu$ is increasing on $[a_0,a_1]$ with $\mu(a_0)=0.18>0$, and
\[
\frac{c_1}{c_2} = 0.55 < \frac{a_1}{a_2}\approx0.778 < \frac{b_1}{b_2} = 4,
\]
so that \eqref{weak_competition} holds. The corresponding steady states are
\[
u^*\approx0.1188,\qquad v^*\approx0.8406,\qquad \frac{a_2}{c_2}=0.9,
\]
and the two thresholds of \eqref{Rstar} are $R_0\approx0.594$ and
$R^*\approx1.097$.

\subsection{The spreading--vanishing dichotomy}

We first fix the shifting speed at $c=0.04$ and vary $h_0$.

\smallskip
\noindent\textbf{Vanishing of $u$.} For $h_0=0.3$ (note $h_0<R_0$) the invasive
species fails to establish itself and dies out; the free boundary converges to
$h_\infty\approx0.91$, which is consistent with the theoretical bound
$h_\infty\le R^*\approx1.097$ of \cref{th:dynamics}(ii). Relieved from the
competitive pressure of $u$, the native species recovers its own carrying
capacity $a_2/c_2=0.9$. See \cref{fig:u_vanishing,fig:u_vanishing_3d,fig:u_vanishing_ht}.

\begin{figure}[htbp]
	\centering
	\begin{subfigure}{0.45\textwidth}
		\centering
		\includegraphics[width=\linewidth]{Images/u_vanishing_a.png}
		\caption{}
		\label{vanishing_u}
	\end{subfigure}
	\hfill
	\begin{subfigure}{0.45\textwidth}
		\centering
		\includegraphics[width=\linewidth]{Images/u_vanishing_b.png}
		\caption{}
		\label{vanishing_v}
	\end{subfigure}
	\caption{Vanishing of the invasive species with $c=0.04$ and $h_0 = 0.3$: (A) heatmap of $u(t,x)$; (B) heatmap of $v(t,x)$.}
	\label{fig:u_vanishing}
\end{figure}

\begin{figure}[htbp]
	\centering
	\begin{subfigure}{0.45\textwidth}
		\centering
		\includegraphics[width=\linewidth]{Images/vanishing_byw_u.png}
		\caption{}
		\label{vanishing_u_3d}
	\end{subfigure}
	\hfill
	\begin{subfigure}{0.45\textwidth}
		\centering
		\includegraphics[width=\linewidth]{Images/vanishing_byw_v.png}
		\caption{}
		\label{vanishing_v_3d}
	\end{subfigure}
	\caption{The same simulation as in \cref{fig:u_vanishing}, shown as surface plots: (A) $u(t,x)$; (B) $v(t,x)$.}
	\label{fig:u_vanishing_3d}
\end{figure}

\begin{figure}[htbp]
	\centering
	\includegraphics[scale=0.8]{Images/u_vanishing_ht.png}
	\caption{The free boundary $h(t)$ when $u$ vanishes ($c = 0.04$, $h_0 = 0.3$). Here $h_\infty\approx0.91\le R^*\approx1.097$.}
	\label{fig:u_vanishing_ht}
\end{figure}

\noindent\textbf{Spreading of $u$.} For $h_0=1$ the invasion succeeds: the range
$[0,h(t)]$ expands indefinitely and the two densities converge to the coexistence
state $(u^*,v^*)\approx(0.12,0.84)$. Ahead of the invasion front the native
species sits at its own carrying capacity $a_2/c_2=0.9$; behind the front it
drops to $v^*\approx0.84$, so that the invasion depresses, but does not
eliminate, the resident population. See \cref{fig:heat_spread,fig:byw}. Note that
$h_0=1<R^*\approx1.097$, so this run also shows that the sufficient condition
$h_0\ge R^*$ of \cref{th:criteria}(i) is not necessary.

\begin{figure}[htbp]
	\centering
	\begin{subfigure}{0.45\textwidth}
		\centering
		\includegraphics[width=\linewidth]{Images/u_spreading_a.png}
		\caption{}
		\label{spreading_heat_u}
	\end{subfigure}
	\hfill
	\begin{subfigure}{0.45\textwidth}
		\centering
		\includegraphics[width=\linewidth]{Images/u_spreading_b.png}
		\caption{}
		\label{spreading_heat_v}
	\end{subfigure}
	\caption{Spreading of the invasive species with $c=0.04$ and $h_0=1$: (A) heatmap of $u(t,x)$; (B) heatmap of $v(t,x)$. The straight interface is the invasion front $x=h(t)$.}
	\label{fig:heat_spread}
\end{figure}

\begin{figure}[htbp]
	\centering
	\begin{subfigure}{0.45\textwidth}
		\centering
		\includegraphics[width=\linewidth]{Images/byw_u.png}
		\caption{}
		\label{u_spreading_a}
	\end{subfigure}
	\hfill
	\begin{subfigure}{0.45\textwidth}
		\centering
		\includegraphics[width=\linewidth]{Images/byw_v.png}
		\caption{}
		\label{u_spreading_b}
	\end{subfigure}
	\caption{The same simulation as in \cref{fig:heat_spread}, shown as surface plots: (A) $u(t,x)$ on the moving domain; (B) $v(t,x)$ on the fixed domain.}
	\label{fig:byw}
\end{figure}

\subsection{The asymptotic spreading speed}

We now fix $h_0=1$, so that spreading occurs, and vary the shifting speed $c$.
\Cref{fig:hovertc004} displays the ratio $h(t)/t$ for $c=0.04$: it stabilises
around $0.0128$, that is, at a value strictly below $c$. This is the regime
$c>c_0$ of \cref{th:spreading_speed}, and the plateau provides a numerical
estimate $c_0\approx0.0128$ of the intrinsic speed. \Cref{fig:compare} and
\cref{table1} compare $h(t)/t$ for several values of $c$: for $c\le0.01$ the
ratio settles near $c$, while for $c\ge0.02$ it settles near the same value
$c_0\approx0.0128$ regardless of $c$. This is in complete agreement with the
formula $\lim_{t\to\infty}h(t)/t=\min\{c,c_0\}$. The slight excess of the tabulated
values over $\min\{c,c_0\}$ for the smaller speeds is a finite-time effect: when
$c<c_0$ one has $h(t)=ct+O(1)$, so that $h(t)/t-c$ decays only like $1/t$.

Finally, running the same code on the single species problem \eqref{single_hom}
(that is, with $c_1=0$ and $A\equiv a_1$) gives $c_0^{\rm sing}\approx0.040$ for
these parameters. The intrinsic speed of the invader is therefore reduced by a
factor of more than three by the mere presence of the native competitor, which
quantifies \cref{lm:c0sing} and \cref{rm:speed}.

\begin{figure}[htbp]
	\centering
	\begin{subfigure}{0.45\textwidth}
		\centering
		\includegraphics[width=\linewidth]{Images/hovertc004_a.png}
		\caption{}
		\label{hovertc004_a}
	\end{subfigure}
	\hfill
	\begin{subfigure}{0.45\textwidth}
		\centering
		\includegraphics[width=\linewidth]{Images/hovertc004_b.png}
		\caption{}
		\label{hovertc004_b}
	\end{subfigure}
	\caption{The ratio $h(t)/t$ for $c = 0.04$: (A) over $t\in[50, 1600]$; (B) a zoom for large $t$.}
	\label{fig:hovertc004}
\end{figure}

\begin{table}[htbp]
	\centering
	\begin{tabular}{c|cccccccc}
		$c$ & $0.007$ & $0.008$ & $0.009$ & $0.01$ & $0.02$ & $0.03$ & $0.04$ & $0.05$\\
		\hline
		$h(t)/t$ at $t\approx1600$ & $0.0074$ & $0.0084$ & $0.0094$ & $0.0103$ & $0.0125$ & $0.0127$ & $0.0128$ & $0.0128$
	\end{tabular}
	\caption{The ratio $h(t)/t$ at $t\approx1600$ for selected values of $c$ ($h_0=1$). The values are consistent with $\min\{c,c_0\}$ and $c_0\approx0.0128$.}
	\label{table1}
\end{table}

\begin{figure}[htbp]
	\centering
	\includegraphics[scale=0.8]{Images/compare.png}
	\caption{Comparison of $h(t)/t$ over $t \in [50,1600]$ for selected values of $c$.}
	\label{fig:compare}
\end{figure}

\section{Concluding remarks}\label{sec_8}

We have determined the long-time dynamics and the exact asymptotic spreading
speed of a free boundary competition system in which an invasive species benefits
from a climate shift while the native competitor occupies a fixed, globally
favourable habitat, in the weak competition regime where the two species may
coexist. The dichotomy of \cref{th:dynamics}, the criteria of
\cref{th:criteria} and the formula $\min\{c,c_0\}$ of
\cref{th:spreading_speed} complete, for the model \eqref{main}, the picture
started in \cite{DLN2026a} for the weak--strong regimes.

Several natural questions remain open. First, as in \cite{HU20205931}, we do not
know whether the threshold $\sigma_0$ of \cref{th:criteria}(ii) can be infinite.
Second, when $c<c_0$ our result gives the speed of the front but not its precise
profile. In the single species case \cite[Theorem 1.2]{HU20205931} shows that
$h(t)-ct\to L_0$ and that the solution converges to the semi-wave
$\psi_{L_0}$, and it is natural to expect the analogous statement here, with
$(\psi_{L_0},\varphi_{L_0})$ given by \cref{th:semi-wave}. Third, the strong
competition case $\frac{b_1}{b_2}<\frac{a_1}{a_2}<\frac{c_1}{c_2}$ is covered
neither by \cite{DLN2026a} nor by the present paper, and this exhausts the
parameter space. In that regime the kinetic system is \emph{bistable}: both
semi-trivial states $(\frac{a_1}{b_1},0)$ and $(0,\frac{a_2}{c_2})$ are locally
stable, one has $b_1c_2-b_2c_1<0$, and the coexistence state \eqref{coexist},
although still positive, is a saddle, so that coexistence is not the generic
outcome. Ecologically this is the priority effect, in which the species that
occupies the habitat first tends to keep it. Two of the tools used here then break down, already for the
homogeneous problem \eqref{main2}. The monotone iteration of \cref{th1:lm2},
which identifies the state behind the front, relies on the existence of a
globally attracting kinetic state and no longer closes. The travelling waves of
a bistable competition system come with a single, possibly negative, speed
rather than with a family parametrised by $c\in[0,c^*)$, so that the
characterisation \eqref{c0} of $c_0$ loses its meaning. The homogeneous problem
has since been settled in paper \cite{DL2026c}, which truncates
the semi-wave problem to a bounded interval, where the sliding method applies
with strict boundary inequalities, and obtains a unique monotone semi-wave for
every speed below the bistable one together with the selection of $c_0(\mu)$.
What does not carry over to the shifting setting is the criterion for the
invader to establish itself: under strong competition the growth rate at the
leading edge is negative, so a small founding population can never invade, and
\cite{DL2026c} replaces the criterion by the domination of a stationary pair on
a bounded interval. The analogue of that pair in a moving favourable region is
not available. Fourth, one may ask what happens when the native
competitor is itself subject to a shifting environment. If its growth rate is
replaced by a function $\tilde a_2(x-ct)$ moving at the \emph{same} speed $c$,
the system remains autonomous in the frame $\xi=x-ct$, so that semi-waves in the
sense of \cref{sec_5} still make sense and the comparison arguments of
\cref{sec_6} are unaffected; what has to be redone is the input from
\cite{MR3986328}, since \cref{th:c*,th:semi-wave0} are stated for a constant
$a_2$, and consequently also the description of the limiting state of $v$ in
\cref{th:dynamics}, which is then an $x$-dependent profile rather than the
constant $\frac{a_2}{c_2}$. The case of two climatic edges moving at different
speeds admits no autonomous moving frame and seems out of reach of the present
techniques. Finally, it
would be interesting to treat the corresponding model with nonlocal diffusion, in
the spirit of \cite{MR4873231}.

\appendix

\section{Proof of \texorpdfstring{\cref{th4}}{the local existence theorem}}\label{sec_app}

We follow the strategy of \cite[Theorem 2.1]{MR3327894}; the details are included
for completeness, since the shifting terms $A(\cdot-ct)$ and $\mu(A(\cdot))$
require some care.

\textit{Step 1: straightening the free boundary.} Let $\zeta\in C^3([0,\infty))$
satisfy
\[
\zeta(y) = 1 \text{ for } |y - h_0| < \tfrac{h_0}{8}, \quad \zeta(y) = 0 \text{ for } |y - h_0| > \tfrac{h_0}{2}, \quad |\zeta'(y)| < \tfrac{6}{h_0}\ \ \forall y,
\]
and set $x=F(t,y):=y+\zeta(y)(h(t)-h_0)$ for $y\ge0$. If
$|h(t)-h_0|\le\frac{h_0}{8}$, then
$\partial_yF=1+\zeta'(y)(h(t)-h_0)\ge1-\frac{6}{h_0}\cdot\frac{h_0}{8}>0$, so
$F(t,\cdot)$ is a diffeomorphism of $[0,\infty)$ mapping $[0,h_0]$ onto
$[0,h(t)]$. Writing
\[
\sqrt{B(h,y)}:=\frac{1}{1+\zeta'(y)(h-h_0)},\quad
C(h,y):=-\frac{\zeta''(y)(h - h_0)}{[1+\zeta'(y)(h-h_0)]^3},\quad
D(h,y):=\frac{\zeta(y)}{1 + \zeta'(y)(h - h_0)},
\]
and $w(t,y):=u(t,F(t,y))$, $z(t,y):=v(t,F(t,y))$, problem \eqref{main} becomes
\begin{equation}\label{eq:transform}
	\begin{cases}
		w_t - d_1 B w_{yy} - (d_1C + h' D) w_y = \big(A(F(t,y) - ct) - b_1w - c_1z \big)w, & t > 0, ~ 0 \le y < h_0,\\
		z_t - d_2 B z_{yy} - (d_2C + h' D) z_y = \left(a_2 -b_2w -c_2z \right)z, & t >0, ~ 0 \le y < \infty,\\
		w_y(t,0) = z_y(t,0) = 0, \quad w(t, y) = 0 \ (y\ge h_0), & t > 0,\\
		h'(t) = -\mu\big(A(h(t) - ct) \big)  w_y(t, h_0), & t > 0,\\
		h(0) = h_0, ~ w(0, y) = u_0(y),\ z(0, y) = v_0(y). &
	\end{cases}
\end{equation}
Set $h^* := -\mu(A(h_0))u_0'(h_0)$, $\Delta_T := [0,T] \times [0,h_0]$ and
$\Delta_T^\infty := [0, T] \times [0, \infty)$. For
$0 < T < \frac{h_0}{8(1+h^*)}$ define $\Gamma_T := W_T \times Z_T \times H_T$,
where
\begin{align*}
	H_T &= \left\{h \in C^1([0, T]) : h(0) = h_0, ~ h'(0) = h^*, ~ \|h' - h^*\|_{C([0,T])} \le 1 \right\}, \\
	W_T &= \left\{ w \in C(\Delta_T^\infty) : w \equiv 0 \text{ for } y \ge h_0, ~ w(0,\cdot) = u_0, ~ \|w - u_0\|_{C(\Delta_T)} \le 1 \right\}, \\
	Z_T &= \left\{ z \in C_b(\Delta_T^\infty) : z(0,\cdot) = v_0, ~ \|z - v_0\|_{L^\infty(\Delta_T^\infty)} \le 1 \right\},
\end{align*}
a complete metric space for
$\mathcal{D}\big((w_1, z_1, h_1), (w_2, z_2, h_2)\big) = \|w_1 - w_2\|_{C(\Delta_T)}
+ \|z_1 - z_2\|_{L^\infty(\Delta_T^\infty)} + \|h_1' - h_2'\|_{C([0,T])}$.

\textit{Step 2: a priori estimates.} For $h_1,h_2 \in H_T$, since
$h_1(0)=h_2(0)$,
\begin{equation}\label{eq:h12}
	\|h_1 - h_2\|_{C([0,T])} \le T \|h_1' - h_2'\|_{C([0,T])}.
\end{equation}
Given $(w,z,h)\in\Gamma_T$ put
$f(t,y):=\big(A(F(t,y)-ct)-b_1w-c_1z\big)w$ and
$g(t,y):=(a_2-b_2w-c_2z)z$. Since $\|w\|_\infty\le M_1':=\|u_0\|_{C([0,h_0])}+1$,
$\|z\|_\infty\le M_2':=\|v_0\|_{L^\infty}+1$ and $A\le a_1$, both $f$ and $g$ are
bounded by constants depending only on the data. By $L^p$ theory and Sobolev
embedding, the linear problem obtained from \eqref{eq:transform} by freezing
$(w,z,h)$ in the right hand side and in the coefficients admits a unique bounded
solution $(\tilde w,\tilde z)$ with
\begin{equation}\label{eq:wz}
	\|\tilde{w} \|_{C^{(1+\alpha)/2,1+\alpha}(\Delta_T)} \le C_1, \qquad \|\tilde{z}\|_{C^{(1+\alpha)/2,1+\alpha}(\Delta_T^\infty)} \le C_1,
\end{equation}
where $C_1$ depends on $\alpha,h_0,c,a_0,a_1,a_2,b_i,c_i,d_i$,
$\|u_0\|_{C^2([0,h_0])}$ and $\|v_0\|_{C^2([0,\infty))}$. Define
\[
\tilde{h}(t) := h_0 - \int_0^t \mu \big( A(h(\tau) - c\tau) \big) \tilde{w}_y(\tau, h_0)\,\mathrm{d}\tau ,
\]
so that $\tilde h(0)=h_0$, $\tilde h'(0)=h^*$ and, $\mu\circ A$ being Lipschitz,
\begin{equation}\label{eq:hprime}
	\|\tilde{h}'\|_{C^{\alpha/2}([0,T])} \le C_2 ,
\end{equation}
with $C_2$ depending on $C_1$, $c$ and the Lipschitz constants of $\mu$ and $A$.

\textit{Step 3: the contraction mapping.} Let
$\mathcal{F}(w,z,h):=(\tilde w,\tilde z,\tilde h)$; a triple in $\Gamma_T$ is a
fixed point of $\mathcal F$ if and only if it solves \eqref{eq:transform}. From
\eqref{eq:wz}--\eqref{eq:hprime},
\begin{align*}
	\|\tilde{h}' - h^*\|_{C([0,T])} &\le C_2 T^{\alpha/2}, &
	\|\tilde{w} - u_0\|_{C(\Delta_T)} &\le C_1 T^{(1+\alpha)/2}, &
	\|\tilde{z} - v_0\|_{L^\infty(\Delta_T^\infty)} &\le C_1 T^{(1+\alpha)/2},
\end{align*}
so $\mathcal F(\Gamma_T)\subset\Gamma_T$ as soon as
$T \le \min\{C_2^{-2/\alpha}, C_1^{-2/(1+\alpha)}\}$.

Let $(w_i,z_i,h_i)\in\Gamma_T$ and
$(\tilde w_i,\tilde z_i,\tilde h_i)=\mathcal F(w_i,z_i,h_i)$, $i=1,2$. By
\eqref{eq:wz}--\eqref{eq:hprime} these satisfy
\[
\|\tilde{w}_i \|_{C^{(1 + \alpha)/2, 1 + \alpha}(\Delta_T)}\le C_1,
\qquad
\|\tilde{z}_i \|_{C^{(1 + \alpha)/2, 1 + \alpha}(\Delta_T^\infty)}\le C_1,
\qquad
\|\tilde{h}_i' \|_{C^{\alpha/2}([0,T])}\le C_2 .
\]
The function
$W := \tilde{w}_1 - \tilde{w}_2$ satisfies
\[
\begin{aligned}
	& W_t - d_1 B(h_2, y) W_{yy} - \big(d_1C(h_2, y) + h_2' D(h_2, y) \big)W_y \\
	& = d_1[B(h_1, y) - B(h_2, y)] \tilde{w}_{1,yy} + d_1[C(h_1, y) - C(h_2, y) ] \tilde{w}_{1,y} \\
	& \quad + [h_1' D(h_1, y) - h_2' D(h_2, y)]\tilde{w}_{1,y} + f_1 - f_2, \qquad t > 0, \; 0 < y < h_0, \\
	& W_y(t,0) = 0, \quad W(t, h_0) = 0, \quad W(0, y) = 0 .
\end{aligned}
\]
Using the $L^p$ estimates for parabolic equations and Sobolev's embedding
theorem, and the fact that $B,C,D$ are Lipschitz in $h$, we obtain
\begin{equation}\label{eq:w12}
	\|\tilde{w}_1 - \tilde{w}_2 \|_{C^{(1 + \alpha)/2, 1 + \alpha}(\Delta_T)}
	\le C_3 \big(\|w_1 - w_2\|_{C(\Delta_T)} +\|z_1 - z_2\|_{L^\infty(\Delta_T^\infty)}   +\|h_1 - h_2\|_{C^1([0,T])} \big),
\end{equation}
and similarly
\begin{equation}\label{eq:z12}
	\|\tilde{z}_1 - \tilde{z}_2\|_{C^{(1+\alpha)/2, 1+\alpha}(\Delta_T^\infty)}
	\le C_4 \big(\|w_1 - w_2\|_{C(\Delta_T)} + \|z_1 - z_2\|_{L^\infty(\Delta_T^\infty)} + \|h_1 - h_2\|_{C^1([0,T])}\big),
\end{equation}
where $C_3,C_4$ depend on $\alpha, h_0, c, C_1, C_2$, the parameters of the
system, the Lipschitz constants of $A$ and $\mu$ and the cutoff $\zeta$. Since
$\tilde h_i'(t)=-\mu\big(A(h_i(t)-ct)\big)\tilde w_{i,y}(t,h_0)$ and
$\mu\circ A$ is Lipschitz, there is $L_{\mu,A}>0$ with
\begin{equation}\label{eq:hp12}
	\|\tilde{h}'_1 - \tilde{h}'_2\|_{C^{\alpha/2}([0,T])}
	\le L_{\mu,A} \big( \|\tilde{w}_{1,y} - \tilde{w}_{2,y}\|_{C^{\alpha/2,0}(\Delta_T)}
	+ \|h_1 - h_2\|_{C([0,T])}\big).
\end{equation}
Combining \eqref{eq:h12}, \eqref{eq:w12}, \eqref{eq:z12} and \eqref{eq:hp12} and
assuming $T\le1$,
\[
\begin{aligned}
	& \|\tilde{w}_1 - \tilde{w}_2\|_{C^{(1+\alpha)/2, 1+\alpha}(\Delta_T)} + \|\tilde{z}_1 - \tilde{z}_2\|_{C^{(1+\alpha)/2, 1+\alpha}(\Delta_T^\infty)} + \|\tilde{h}'_1 - \tilde{h}'_2\|_{C^{\alpha/2}([0,T])} \\
	& \quad \le C_5 \big(\|w_1 - w_2\|_{C(\Delta_T)} + \|z_1 - z_2\|_{L^\infty(\Delta_T^\infty)} + \|h'_1 - h'_2\|_{C([0,T])}\big),
\end{aligned}
\]
with $C_5$ depending on $C_3, C_4$ and $L_{\mu,A}$. Hence, choosing
\[
T := \min \left\{ 1, \ \left(\tfrac{1}{2C_5}\right)^{2/\alpha}, \ C_2^{-2/\alpha}, \ C_1^{-2/(1+\alpha)}, \ \tfrac{h_0}{8(1+h^*)} \right\},
\]
we get
\[
\begin{aligned}
	& \|\tilde{w}_1 - \tilde{w}_2\|_{C(\Delta_T)} + \|\tilde{z}_1 - \tilde{z}_2\|_{L^\infty(\Delta_T^\infty)} + \|\tilde{h}'_1 - \tilde{h}'_2\|_{C([0,T])}\\
	&\quad \le C_5 T^{\alpha/2} \big(\|w_1 - w_2\|_{C(\Delta_T)} + \|z_1 - z_2\|_{L^\infty(\Delta_T^\infty)} + \|h'_1 - h'_2\|_{C([0,T])}\big) \\
	&\quad \le \tfrac{1}{2} \big(\|w_1 - w_2\|_{C(\Delta_T)} + \|z_1 - z_2\|_{L^\infty(\Delta_T^\infty)} + \|h'_1 - h'_2\|_{C([0,T])}\big).
\end{aligned}
\]
Thus $\mathcal F$ is a contraction on $\Gamma_T$ and has a unique fixed point
$(w,z,h)$. Undoing the change of variables, $(u,v,h)$ solves \eqref{main}, and
Schauder estimates give the additional regularity
\[
h \in C^{1+\alpha/2}([0,T]), \quad u \in C^{1+\alpha/2, 2+\alpha}(D_T), \quad v \in C^{1+\alpha/2, 2+\alpha}((0,T] \times (0, +\infty)),
\]
so that $(u,v,h)$ is a classical solution. Since all the constants above depend
continuously on the data and on $c$, the solution depends continuously on
$(u_0,v_0,h_0,c,\mu)$, a fact used in \cref{sec_4} and \cref{sec_6}. \qed

%\section*{Declarations}
%
%\textbf{Competing interests:} The authors declare that they have no conflict of interest.
%
%\textbf{Funding:} P. Le is supported by University of Economics and Law, Vietnam National University, Ho Chi Minh City, Vietnam; P.V. Dinh is supported by FPT University, Ho Chi Minh City.
%
%\textbf{Availability of data and materials:} not applicable.
%
%\section*{Acknowledgments}
%We are grateful to the anonymous reviewers for their careful reading and helpful suggestions.

\bibliographystyle{abbrv}
\bibliography{ref}

\end{document}
